\documentclass[../Main.tex]{subfiles}


\begin{document}

\chapter{Independence, Conditioning and Asymptotic Theory}


\section{ Measure product and Independence} \label{sec:measure_product_and_independence}
The notion of independence is one of the first genuinely new concepts introduced in probability theory, in the sense that it does not arise as a straightforward adaptation of a corresponding notion from measure theory. While the independence of two events or two random variables admits a natural intuitive interpretation, the most fundamental and structurally relevant notion is that of independence between two or more \sfields.

A central result shows that two random variables are independent if and only if the joint distribution of the pair coincides with the product of their marginal distributions. Combined with Fubini's theorem, which we present first, this characterization yields several equivalent and practical formulations of independence.

We then establish the classical Borel-Cantelli lemma and illustrate its strength through an application to the dyadic expansion of a randomly chosen real number, leading to several striking probabilistic properties.

\subsection{Probabilities measures product}
Let $(\measspace[1])$ and $(\measspace[2])$ be two probabilistic spaces. The product \sfield $\A_1\otimes\A_2$ is defined as the \sfield on $\Omega_1\times\Omega_2$ given by
\[ \A_1\otimes\A_2 = \sigma\bigl(A\times B : A\in\A_1,\ B\in\A_2\bigr). \]

Sets of the form $A\times B$, with $A\in\A_1$ and $B\in\A_2$, are called \emph{measurable rectangles}. 

\textbf{remark :} The definition of the product \sfield extends naturally to a finite collection of probabilistic spaces $(\Omega_1,\A_1),\dots,(\Omega_n,\A_n)$ by setting
\[ \A_1\otimes\cdots\otimes\A_n = \sigma(\A_1\times\cdots\times\A_n). \]
The expected associativity properties hold; for instance, if $n=3$,
\[ (\A_1\otimes\A_2)\otimes\A_3 = \A_1\otimes(\A_2\otimes\A_3) = \A_1\otimes\A_2\otimes\A_3. \]

Let $C\subset \Omega_1\times\Omega_2$ and let $\omega_1\in\Omega_1$. We define the \emph{section} of $C$ at $\omega_1$ by
\[ C_{\omega_1} = \{\omega_2\in\Omega_2 : (\omega_1,\omega_2)\in C\}. \]
Similarly, for $\omega_2\in\Omega_2$, define
\[ C^{\omega_2} = \{\omega_1\in\Omega_1 : (\omega_1,\omega_2)\in C\}. \]

If $X$ is a random variable on $\Omega_1\times\Omega_2$ and $\omega_1\in\Omega_1$, we denote by $X_{\omega_1}$ the random variable on $\Omega_2$ defined by
\[ X_{\omega_1}(\omega_2)=X(\omega_1,\omega_2). \]
Similarly, for $\omega_2\in\Omega_2$, we define $X^{\omega_2}(\omega_1)=X(\omega_1,\omega_2)$.

\proppr{product_sigma_field_sections}{
\begin{enumerate}
\item If $C\in\A_1\otimes\A_2$, then for every
$\omega_1\in\Omega_1$, the section $C_{\omega_1}$ belongs to $\A_2$,
and for every $\omega_2\in\Omega_2$, the section $C^{\omega_2}$ belongs to
$\A_1$.

\item Let $X$ be a random variable on $\Omega_1\times\Omega_2$ with respect to
$\A_1\otimes\A_2$. Then, for every $\omega_1\in\Omega_1$,
$X_{\omega_1}$ is $\A_2$-measurable, and for every
$\omega_2\in\Omega_2$, $X^{\omega_2}$ is $\A_1$-measurable.
\end{enumerate}
}{
\begin{enumerate}
    \item Fix $\omega_1 \in \Omega_1$ and  Consider the class 
    $$ \G=\{C\in\A_1\otimes\A_2: C_{\omega_1}\in \A_2 \}. $$
    Let $ \R=\{A\times B:\;A\in\A_1,\;B\in\A_2\}. $
    One checks that: 
    \begin{itemize}
        \item $\R$ is an algebra generating $\A_1\otimes\A_2$; 
        \item $\R\subset\G$ (if $C = A \times B$, we have $C_{\omega_1} = B$ or $C_{\omega_1} = \vide$ according as $\omega_1 \in A$ or $\omega_1 \notin A$) ; 
        \item $\G$ is a monotone class.
    \end{itemize}
    Hence, by the Monotone Class Theorem, 
    $$ \A_1\otimes\A_2=\sigma(\R)\subset\G, $$ 
    proving that $C_{\omega_1}\in\A_2$ for every $ C \in\A_1\otimes \A_2$. The property $C^{\omega_2} \in \A_1$ is derived in a similar manner.

    \item For every measurable subset $D$ of $G$,
    \begin{align*}
    X_{\omega_1}^{-1}(D)
    &= \{ \omega_2 \in \Omega_2 : X_{\omega_1}(\omega_2) \in D \} \\
    &= \{ \omega_2 \in \Omega_2 : (\omega_1,\omega_2) \in X^{-1}(D) \} \\
    &= (X^{-1}(D))_{\omega_1}.
    \end{align*}

    which belongs to $\A_2$ by part (1).
\end{enumerate}

}

\thmpr{Existence of product measure}{product_measure_existence}{
Let $\PP_1$ and $\PP_2$ be two probability measures on $(\measspace[1])$ and $(\measspace[2])$, respectively.
\begin{enumerate}
\item There exists a unique measure $\PP_1\otimes\PP_2$ on $(\Omega_1\times\Omega_2,\A_1\otimes\A_2)$ such that
\[ (\PP_1\otimes\PP_2)(A\times B) = \PP_1(A)\,\PP_2(B), \qquad A\in\A_1,\ B\in\A_2, \]


\item For every $C\in\A_1\otimes\A_2$, the mappings
\[ \omega_1\mapsto \PP_2(C_{\omega_1}) \quad\text{and}\quad \omega_2\mapsto \PP_1(C^{\omega_2}) \]
are measurable (ie random variables), and
\[ (\PP_1\otimes\PP_2)(C) = \int_{\Omega_1}\PP_2(C_{\omega_1})\,d\PP_1(\omega_1) = \int_{\Omega_2}\PP_1(C^{\omega_2})\,d\PP_2(\omega_2). \]
\end{enumerate}
}{
\begin{enumerate}
    \item \textbf{Uniqueness.}
    Left as an exercise. Apply the Monotone Class Theorem (Theorem~\ref{thm:monotone_class_theorem}) to
    $$
    \C=\{D\in\A_1\otimes\A_2:\; p(D)=q(D)\},
    $$
    where $p$ and $q$ are probability measures satisfying
    $$
    p(A\times B)=q(A\times B)=\PP_1(A)\PP_2(B),
    \qquad A\in\A_1,\;B\in\A_2.
    $$
    (See Figure~\ref{fig:monotone_class_strategy} for the general proof strategy.)

    \textbf{Existence.}
    Motivated by the last identity stated in the theorem, define
    $$
    \PP(C):=\int_{\Omega_1}\PP_2(C_{\omega_1})\,d\PP_1(\omega_1), \qquad C\in\A_1\otimes\A_2.
    $$
    We verify that this mapping is well defined and defines a probability measure satisfying
    $$
    \PP(A\times B)=\PP_1(A)\PP_2(B), \qquad A\in\A_1,\;B\in\A_2.
    $$
    \begin{itemize}

        \item \textit{Well-definedness.}
        By Proposition~\ref{prop:product_sigma_field_sections}, the section $C_{\omega_1}$ belongs to $\A_2$, so $\PP_2(C_{\omega_1})$ is well defined. It remains to prove that
        $$
        \phi_C:\omega_1\longmapsto\PP_2(C_{\omega_1})
        $$
        is $\A_1$-measurable.

        Consider the class
        $$
        \G=\{C\in\A_1\otimes\A_2:\phi_C
        \text{ is }\A_1\text{-measurable}\}.
        $$
        ANd use the Monotone Class Theorem to prove :
        $$ \A_1\otimes\A_2=\sigma(\R)\subset\G, $$
        proving that $\phi_C$ is measurable for every $C\in\A_1\otimes\A_2$.

        \item \textit{$\PP$ is a measure.}
        Let $(C_n)_{n\ge 1}$ be pairwise disjoint subsets of $\A_1\otimes\A_2$. Since the sections $(C_n)_{\omega_1}$ are pairwise disjoint for every $\omega_1$,
        $$
        \begin{aligned}
        \PP\!\left(\bigcup_{n=1}^\infty C_n\right)
        &=\int_{\Omega_1}\PP_2\!\left(\left(\bigcup_{n=1}^\infty C_n\right)_{\omega_1}\right)\,d\PP_1(\omega_1)\\
        &=\int_{\Omega_1}\sum_{n=1}^\infty\PP_2((C_n)_{\omega_1})\,d\PP_1(\omega_1)\\
        &=\sum_{n=1}^\infty\int_{\Omega_1}\PP_2((C_n)_{\omega_1})\,d\PP_1(\omega_1)\\
        &=\sum_{n=1}^\infty\PP(C_n),
        \end{aligned}
        $$
        where the third equality follows from the Monotone Convergence Theorem. Thus $\PP$ is countably additive.

        Finally, for every $A\in\A_1$ and $B\in\A_2$,
        $$
        \PP(A\times B)
        =\int_{\Omega_1}\PP_2((A\times B)_{\omega_1})\,d\PP_1(\omega_1)
        =\PP_1(A)\PP_2(B).
        $$
        Therefore, $\PP$ is a probability measure satisfying the required property.
    \end{itemize}
\end{enumerate}

The final identity of the theorem is established similarly. Indeed, the corresponding integrand is well defined by the same measurability argument used above.    
}


\textbf{Remarks:} 
\begin{enumerate}
\item Suppose now that we consider an arbitrary finite number of $\sigma$-finite measures $\PP_1,\ldots,\PP_n$, defined on $(E_1,\A_1),\ldots,(E_n,\A_n)$ respectively. We can then define the product measure
\[ \PP_1 \otimes \cdots \otimes \PP_n \]
on
\[ (E_1 \times \cdots \times E_n,\ \A_1 \otimes \cdots \otimes \A_n) \]
by setting
\[ \PP_1 \otimes \cdots \otimes \PP_n = \PP_1 \otimes (\PP_2 \otimes (\cdots \otimes \PP_n)). \]
The way we insert parentheses is in fact unimportant since $\PP_1 \otimes \cdots \otimes \PP_n$ is characterized by its values on (measurable) rectangles,
\[ \PP_1 \otimes \cdots \otimes \PP_n(A_1 \times \cdots \times A_n) = \PP_1(A_1)\cdots\PP_n(A_n). \]
\end{enumerate}


\thmpr{Fubini-Tonelli Theorem}{Fubini_Tonelli}{
Let $(\measspace[1])$ and $(\measspace[2])$ be probabilistic spaces, and equip $\Omega_1\times\Omega_2$ with the product \sfield $\A_1\otimes\A_2$ and Let $\PP_1$ and $\PP_2$ be two probability measures on $(\measspace[1])$ and $(\measspace[2])$, respectively, and let $X:\Omega_1\times\Omega_2\to[0,\infty]$ be a positive random variable.
\begin{enumerate}
\item The functions
\[\Omega_1\ni \omega_1 \longmapsto \int_{\Omega_2} X(\omega_1,\omega_2)\, d\PP_2(\omega_2), \qquad \Omega_2\ni \omega_2 \longmapsto \int_{\Omega_1} X(\omega_1,\omega_2)\, d\PP_1(\omega_1), \]
with values in $[0,\infty]$, are $\A_1$-measurable and $\A_2$-measurable, respectively (ie. they are random variables).

\item We have :
\begin{align*}
\int_{\Omega_1 \times \Omega_2} X \, d(\PP_1 \otimes \PP_2)
&= \int_{\Omega_1} \left( \int_{\Omega_2} X(\omega_1,\omega_2)\, d\PP_2(\omega_2)\right)d\PP_1(\omega_1) \\
&= \int_{\Omega_2}\left(\int_{\Omega_1}X(\omega_1,\omega_2)\, d\PP_1(\omega_1)\right)d\PP_2(\omega_2).
\end{align*}

Or more simply :
\[ \EE_{\PP_1 \otimes \PP_2}[X] = \EE_{\PP_1}\Big[
      \EE_{\PP_2}\big[ X(\omega_1,\cdot) \big] \Big] = \EE_{\PP_2}
   \Big[ \EE_{\PP_1}\big[ X(\cdot,\omega_2) \big]\Big].
\]

\end{enumerate}
}{
We give the proof for the special case $X=\1_C$ for some $C\in\A_1\otimes\A_2$. The general case can be deduced from this special case by applying the monotone convergence theorem to the sequence of random variables $X_n \conv{n} X$. \\
In fact, let $X=\1_C$ for some $C\in\A_1\otimes\A_2$ the functions in the theorem are reduced to : 
\[\Omega_1\ni \omega_1 \longmapsto \int_{\Omega_2} \1_{C_{\omega_1}}(\omega_2), d\PP_2(\omega_2), \qquad \Omega_2\ni \omega_2 \longmapsto \int_{\Omega_1} \1_{C^{\omega_2}}(\omega_1)\, d\PP_1(\omega_1), \]
Which are exactly : 
\[ \omega_1\mapsto \PP_2(C_{\omega_1}) \quad\text{and}\quad \omega_2\mapsto \PP_1(C^{\omega_2}) \]
By the second point of Theorem \ref{thm:product_measure_existence}, these functions are measurable, and we have
\[ (\PP_1\otimes\PP_2)(C) = \int_{\Omega_1}\PP_2(C_{\omega_1})\,d\PP_1(\omega_1) = \int_{\Omega_2}\PP_1(C^{\omega_2})\,d\PP_2(\omega_2). \]
Which is :
\[ \EE_{\PP_1 \otimes \PP_2}[X] = \EE_{\PP_1}\Big[
      \EE_{\PP_2}\big[ X(\omega_1,\cdot) \big] \Big] = \EE_{\PP_2}
   \Big[ \EE_{\PP_1}\big[ X(\cdot,\omega_2) \big]\Big].
\]
}

\thmpr{Fubini-Lebesgue Theorem}{Fubini_Lebesgue}{
Let $X\in L^1(\Omega_1\times\Omega_2,\A_1\otimes\A_2, \PP_1\otimes\PP_2)$. Then:
\begin{enumerate}
\item For $\PP_1$-almost every $\omega_1\in\Omega_1$, the random variable $\omega_2\mapsto X(\omega_1,\omega_2)$ belongs to $L^1(\probspace[2])$, and for $\PP_2$-almost every $\omega_2\in\Omega_2$, the random variable $\omega_1\mapsto X(\omega_1,\omega_2)$ belongs to $L^1(\probspace[1])$.

\item The functions
\[ \omega_1 \longmapsto \int_{\Omega_2} X(\omega_1,\omega_2)\, d\PP_2(\omega_2), \qquad \omega_2 \longmapsto \int_{\Omega_1} X(\omega_1,\omega_2)\, d\PP_1(\omega_1), \]
belong to $L^1(\probspace[1])$ and $L^1(\probspace[2])$, respectively.

\item We have :
\begin{align*}
\int_{\Omega_1 \times \Omega_2} X \, d(\PP_1 \otimes \PP_2)
&= \int_{\Omega_1} \left( \int_{\Omega_2} X(\omega_1,\omega_2)\, d\PP_2(\omega_2)\right)d\PP_1(\omega_1) \\
&= \int_{\Omega_2}\left(\int_{\Omega_1}X(\omega_1,\omega_2)\, d\PP_1(\omega_1)\right)d\PP_2(\omega_2).
\end{align*}

Or more simply :
\[ \EE_{\PP_1 \otimes \PP_2}[X] = \EE_{\PP_1}\Big[
      \EE_{\PP_2}\big[ X(\omega_1,\cdot) \big] \Big] = \EE_{\PP_2}
   \Big[ \EE_{\PP_1}\big[ X(\cdot,\omega_2) \big]\Big].
\]

\end{enumerate}
}{
\begin{enumerate}
\item  By theorem \ref{thm:Fubini_Tonelli} applied to $|X|$, we have
\[ \int_{\Omega_1} \left( \int_{\Omega_2} |X(\omega_1,\omega_2)| \, d\PP_2(\omega_2) \right)d\PP_1(\omega_1) = \int |X| \, d\PP_1 \otimes \PP_2 < \infty. \]

Thanks to propertry (5) following definition of expectation \ref{defn:Expectation}, this implies that
\[ \int_{\Omega_2} |X(\omega_1,\omega_2)| \, d\PP_2(\omega_2) < \infty \quad \text{for } \PP_1\text{-a.e. } \omega_1. \]
Thus, the function $\omega_1 \mapsto X(\omega_1,\omega_2)$ (which is measurable by proposition \ref{prop:product_sigma_field_sections}) belongs to $L^1(\Omega_2,\A_2,\PP_2)$ except on the $\A_1$-measurable set
\[ N := \left\{ \omega_1 \in \Omega_1 : \int_{\Omega_2} |X(\omega_1,\omega_2)| \, d\PP_2(\omega_2) = \infty \right\}, \]
which has zero $\PP_1$-measure. The same argument applies to the function $\omega_2 \mapsto X(\omega_1,\omega_2)$.

\medskip

\item  If $\omega_1 \in N^c$, $\int_{\Omega_2} X(\omega_1,\omega_2)\,d\PP_2(\omega_2)$ is well defined since the function is only defined (thanks to (1)) outside a measurable set of values of $\omega_1$ of zero $\PP_1$-measure. We can always agree that that $\int_{\Omega_2} X(\omega_1,\omega_2)\,d\PP_2(\omega_2)=0$ if $\omega_1 \in N$. Then the function
\begin{align*}
\omega_1 \mapsto \int_{\Omega_2} X(\omega_1,\omega_2)\,d\PP_2(\omega_2)
&= \1_{N^c}(\omega_1)\int_{\Omega_2} X^+(\omega_1,\omega_2)\,d\PP_2(\omega_2) \\
&\quad - \1_{N^c}(\omega_1)\int_{\Omega_2} X^-(\omega_1,\omega_2)\,d\PP_2(\omega_2).
\end{align*}

is $\A_1$-measurable by \ref{thm:Fubini_Tonelli} (1), and we have also, thanks to part (2) of the same theorem,
\begin{align*}
\int_{\Omega_1} \left| \int_{\Omega_2} X(\omega_1,\omega_2)\,d\PP_2(\omega_2) \right| d\PP_1(\omega_1)
&\leq \int_{\Omega_1} \left( \int_{\Omega_2} |X(\omega_1,\omega_2)|\,d\PP_2(\omega_2) \right)d\PP_1(\omega_1) \\
&= \int |X| \, d\PP_1 \otimes \PP_2 \\
&< \infty.
\end{align*}

This proves that the function $\omega_1 \mapsto \int_{\Omega_2} X(\omega_1,\omega_2)\,d\PP_2(\omega_2)$ is in $L^1(\Omega_1,\A_1,\PP_1)$.
The same argument applies to the function $\omega_2 \mapsto \int_{\Omega_1} X(\omega_1,\omega_2)\,d\PP_1(\omega_1)$.

\medskip

\item Theorem \ref{thm:Fubini_Tonelli} gives the two equalities
\begin{align*}
\int_{\Omega_1} \left( \int_{\Omega_2} X^+(\omega_1,\omega_2)\,d\PP_2(\omega_2) \right)d\PP_1(\omega_1)
&= \int_{\Omega_1 \times \Omega_2} X^+ \, d\PP_1 \otimes \PP_2, \\
\int_{\Omega_1} \left( \int_{\Omega_2} X^-(\omega_1,\omega_2)\,d\PP_2(\omega_2) \right)d\PP_1(\omega_1)
&= \int_{\Omega_1 \times \Omega_2} X^- \, d\PP_1 \otimes \PP_2.
\end{align*}

We can replace the integral over $\Omega_1$ by an integral over $N^c$ in the left-hand sides
of these formulas. Then we just have to subtract the second equality from the first one
(using part (2)), in order to arrive at the first equality in (3), and the second equality
is derived in a similar manner.
\end{enumerate}
}

\subsection{Convolution}

Throughout this section, we consider the measure space $(\RR^d, \B(\RR^d), \lambda)$. Let $f$ and $g$ be two real measurable functions on $\RR^d$. For $x \in \RR^d$, the convolution
\[ (f * g)(x) = \int_{\RR^d} f(x-y)g(y)\,dy \]
is well defined provided
\[ \int_{\RR^d} |f(x-y)g(y)|\,dy < \infty. \]

In that case, the fact that Lebesgue measure on $\RR^d$ is invariant under translations and under the symmetry $y \mapsto -y$ shows that $(g * f)(x)$ is also well defined and $(g * f)(x) = (f * g)(x)$.

\proppr{conv_Lp_inequ}{
Let $f \in L^1(\RR^d,\B(\RR^d),\lambda)$ and $g \in L^p(\RR^d,\B(\RR^d),\lambda)$, where $p \in [1,\infty)$. \\ 
Then, for $\lambda$ a.e. $x \in \RR^d$\footnote{ for all $x$ except a set of null measure}, the convolution $(f * g)(x)$ is well defined.\\
Moreover $f * g \in L^p(\RR^d,\B(\RR^d),\lambda)$ and : 
    \[\|f * g\|_p \leq \|f\|_1 \|g\|_p\]
}{
Again, we may agree that, on the set of zero Lebesgue measure where $f * g$ is not well defined, we take $f * g(x) = 0$.

We leave aside the case $p = \infty$, which is easier (and where stronger results are proved in Proposition \ref{prop:conv_Lpq_inequ} below). We assume that $\|f\|_1 > 0$, since otherwise the results are trivial. Since the measure with density $\|f\|_1^{-1}|f(x)|$ with respect to $\lambda$ is a probability measure, Jensen’s inequality gives, for every $x \in \RR^d$,
\begin{align*}
\left(\int_{\RR^d} |g(x-y)||f(y)|\,dy \right)^p
    &= \|f\|_1^p \left(\int_{\RR^d} |g(x-y)| \frac{|f(y)|}{\|f\|_1} dy \right)^p \\
    &\leq \|f\|_1^{p-1} \int_{\RR^d} |g(x-y)|^p |f(y)|\,dy.
\end{align*}

Using Fubini-Tonelli Theorem \ref{thm:Fubini_Tonelli}, we have then :
\begin{align*}
\int_{\RR^d} \left(\int_{\RR^d} |g(x-y)||f(y)|\,dy \right)^p dx
    &\leq \|f\|_1^{p-1} \int_{\RR^d} \left(\int_{\RR^d} |g(x-y)|^p |f(y)|\,dy \right) dx \\
    &= \|f\|_1^{p-1} \int_{\RR^d} \left(\int_{\RR^d} |g(y)|^p |f(x-y)|\,dy \right) dx \\
    &= \|f\|_1^{p-1} \int_{\RR^d} |g(y)|^p \left(\int_{\RR^d} |f(x-y)| dx \right) dy \\
    &= \|f\|_1^p \|g\|_p^p < \infty,
\end{align*}
which shows that
\[
\int_{\RR^d} |f(x-t)||g(t)|dt < \infty, \qquad \lambda(dx)\ \text{a.e.} \]


and gives the first assertion. The second one also follows since the previous calculation gives
\[
\int_{\RR^d} |f * g(x)|^p dx
\leq \int_{\RR^d} \left(\int_{\RR^d} |g(x-y)||f(y)|\,dy \right)^p dx
\leq \|f\|_1^p \|g\|_p^p.
\]
}

The next proposition gives another important case where the convolution $f * g$ is well defined and enjoys nice continuity properties.

\proppr{conv_Lpq_inequ}{
Let $p \in [1,\infty)$ and $q \in (1,\infty)$ be conjugate exponents $\left(\frac{1}{p} + \frac{1}{q} = 1\right)$. Let $f \in L^p(\RR^d,\B(\RR^d),\lambda)$ and $g \in L^q(\RR^d,\B(\RR^d),\lambda)$. Then, the convolution $f * g(x)$ is well defined and $|f * g(x)| \leq \|f\|_p \|g\|_q$, for every $x \in \RR^d$. Furthermore, the function $f * g$ is uniformly continuous on $\RR^d$.
}{
By the Hölder inequality,
\[
\int_{\RR^d} |f(x-y)g(y)|dy
\leq \left(\int |f(x-y)|^p dy \right)^{1/p} \|g\|_q
= \|f\|_p \|g\|_q.
\]

This gives the first assertion and also proves that $f * g$ is bounded above by $\|f\|_p \|g\|_q$. To get the property of uniform continuity, we use the Lemma \ref{lem:continuity_translation} it is easy to complete the proof of the proposition. Indeed, for every $x,x' \in \RR^d$,
\begin{align*}
|f * g(x) - f * g(x')|
&\leq \int |f(x-y) - f(x'-y)||g(y)|dy \\
&\leq \|g\|_q \left(\int |f(x-y) - f(x'-y)|^p dy \right)^{1/p} \\
&= \|g\|_q \|f \circ \sigma_{-x} - f \circ \sigma_{-x'}\|_p
\end{align*}
and we use Lemma \ref{lem:continuity_translation} to say that $\|f \circ \sigma_{-x} - f \circ \sigma_{-x'}\|_p$ tends to $0$ when $x-x' \to 0$.
}

\lempr{Continuity of translation}{continuity_translation}{
Set $\sigma_x(y) = y - x$ for every $x,y \in \RR^d$. Let $p \in [1,\infty)$ and $f \in L^p(\RR^d,\B(\RR^d),\lambda)$. Then the mapping $x \mapsto f \circ \sigma_x$ is uniformly continuous from $\RR^d$ into $L^p(\RR^d,\B(\RR^d),\lambda)$.
}{
The proof of this lemma is omitted, but could be found in the book \parencite[lem~5.7]{Probability}
%  Suppose first that $f$ belongs to the space $C_c(\RR^d)$ of all real continuous functions with compact support on $\RR^d$. Then,
% \[
% \|f \circ \sigma_x - f \circ \sigma_y\|_p^p
% = \int |f(z-x) - f(z-y)|^p dz
% = \int |f(z) - f(z-(y-x))|^p dz
% \]
% which tends to $0$ when $y-x \to 0$ by an application of the dominated convergence theorem. In the general case, Theorem 4.8 (3) allows us to find a sequence $(f_n)_{n \in \NN}$ in $C_c(\RR^d)$ that converges to $f$ in $L^p(\RR^d,\B(\RR^d),\lambda)$. Then, for every $x,y \in \RR^d$,
% \begin{align*}
% \|f \circ \sigma_x - f \circ \sigma_y\|_p
% &\leq \|f \circ \sigma_x - f_n \circ \sigma_x\|_p
% + \|f_n \circ \sigma_x - f_n \circ \sigma_y\|_p
% + \|f_n \circ \sigma_y - f \circ \sigma_y\|_p \\
% &= 2\|f - f_n\|_p + \|f_n \circ \sigma_x - f_n \circ \sigma_y\|_p.
% \end{align*}

% Fix $\varepsilon > 0$ and first choose $n$ large enough so that $\|f - f_n\|_p < \varepsilon/4$. Since $f_n \in C_c(\RR^d)$, we can then choose $\delta > 0$ such that $\|f_n \circ \sigma_x - f_n \circ \sigma_y\|_p \leq \varepsilon/2$ if $|x-y| < \delta$. The preceding display then shows that $\|f \circ \sigma_x - f \circ \sigma_y\|_p \leq \varepsilon$ if $|x-y| < \delta$. This completes the proof.
}


\paragraph{Approximations of the Dirac Measure}
\defnr{Approximation of the Dirac Measure}{dirac_aprox}{
Let us say that a sequence $(\varphi_n)_{n \in \NN}$ of nonnegative continuous functions on $\RR^d$ is an approximation of the Dirac measure $\delta_0$ if:
\begin{enumerate}
\item For every $n \in \NN$,
\[ \int_{\RR^d} \varphi_n(x)\,dx = 1. \]

\item For every $\delta > 0$,
\[ \lim_{n \to \infty} \int_{\{|x|>\delta\}} \varphi_n(x)\,dx = 0. \]
\end{enumerate}
}

It is easy to construct approximations of $\delta_0$. If $\varphi$ is any nonnegative continuous function on $\RR^d$ such that $\int \varphi(x)\,dx = 1$, we can set $\varphi_n(x) = n^d \varphi(nx)$ for every $x \in \RR^d$ and every $n \in \NN$. Note that, if we start from a function $\varphi$ with compact support, all functions $\varphi_n$ will vanish outside the same closed ball.

\proppr{dirac_limit}{
Let $(\varphi_n)_{n \in \NN}$ be an approximation of the Dirac measure $\delta_0$.
\begin{enumerate}
\item For any bounded continuous function $f : \RR^d \longrightarrow \RR$, the sequence $\varphi_n * f$ converges to $f$ as $n \to \infty$, uniformly on every compact subset of $\RR^d$.
\item Let $f \in L^p(\RR^d,\B(\RR^d),\lambda)$, where $p \in [1,\infty)$. Then $\varphi_n * f \to f$ in $L^p$.
\end{enumerate}
}{
\begin{enumerate}
\item First note that $\varphi_n * f$ is well defined since $\varphi_n$ is integrable and $f$ is bounded. Then, we have, for every $\delta > 0$,
\begin{align*}
\varphi_n * f(x) - f(x)
&= \int_{|y|\leq \delta} (f(x-y) - f(x))\varphi_n(y)\,dy \\
&\quad + \int_{|y|>\delta} (f(x-y) - f(x))\varphi_n(y)\,dy. \tag{i}
\end{align*}

Fix $\varepsilon > 0$ and a compact subset $H$ of $\RR^d$. Since $f$ is uniformly continuous on any compact subset of $\RR^d$, we can choose $\delta > 0$ such that $|f(x-y)-f(x)| \leq \varepsilon$ for every $x \in H$ and $y \in \RR^d$ such that $|y| \leq \delta$. Recalling (1) in Definition \ref{defn:dirac_aprox}, we get that the first term in the right-hand side of (i) is smaller than $\varepsilon$ in absolute value, for every $x \in H$ and $n \in \NN$. Furthermore, (2) in Definition \ref{defn:dirac_aprox} and the boundedness of $f$ show that the second term in the right-hand side of (i) tends to $0$ when $n \to \infty$, uniformly when $x$ varies in $\RR^d$. It follows that we have $|\varphi_n * f(x) - f(x)| < 2\varepsilon$ for every $x \in H$, as soon as $n$ is sufficiently large.

\item The fact that $\varphi_n * f$ is well defined a.e. and belongs to $L^p$ follows from Proposition \ref{prop:conv_Lp_inequ} since $\varphi_n$ is integrable. We then observe that
\begin{align*}
\int |\varphi_n * f(x) - f(x)|^p dx
    &\leq \int \left( \int \varphi_n(y)|f(x-y)-f(x)|\,dy \right)^p dx \\
    &\leq \int \left( \int \varphi_n(y)|f(x-y)-f(x)|^p dy \right) dx \\
    &= \int \left( \int \varphi_n(y)|f(x-y)-f(x)|^p dx \right) dy \\
    &= \int \varphi_n(y)\left( \int |f(x-y)-f(x)|^p dx \right) dy
\end{align*}

where the second inequality is a consequence of Jensen's inequality (note that $\varphi_n(y)dy$ is a probability measure), and the next equality uses Fubini-Tonelli Theorem \ref{thm:Fubini_Tonelli}. For every $y \in \RR^d$, set
\[ h(y) = \int |f(x-y)-f(x)|^p dx = \|f \circ \sigma_y - f\|_p^p, \]
with the notation of Lemma \ref{lem:continuity_translation}. The function $h$ is bounded and $h(y) \to 0$ as $y \to 0$ by Lemma \ref{lem:continuity_translation}. It then immediately follows from (1) and (2) in Definition \ref{defn:dirac_aprox} that
\[ \lim_{n \to \infty} \int \varphi_n(y)h(y)\,dy = 0. \]
This completes the proof.
\end{enumerate}

}

\paragraph{Application}
In dimension $d = 1$, we can take
\[
\varphi_n(x) = c_n(1 - x^2)^n \1_{\{|x|\leq 1\}}
\]
where the constant $c_n$ is chosen so that $\int \varphi_n(x)dx = 1$. Then let $[a,b]$ be a compact subinterval of $(0,1)$, and let $f$ be a continuous function on $[a,b]$. It is easy to extend $f$ to a continuous function on $\RR$ that vanishes outside $[0,1]$ (for instance, take $f$ to be affine on both intervals $[0,a]$ and $[b,1]$). Then the preceding proposition shows that
\[
\varphi_n * f(x)
= c_n \int (1-(x-y)^2)^n \1_{\{|x-y|\leq 1\}} f(y)\,dy
\longrightarrow f(x)
\]
uniformly on $[a,b]$. If $x \in [a,b]$, we can remove the indicator function $\1_{\{|x-y|\leq 1\}}$ in the last display, since the conditions $x \in [a,b]$ and $|x-y|>1$ force $f(y)=0$. It now follows that $f$ is a uniform limit of polynomials on the interval $[a,b]$ (this is a special case of the Stone-Weierstrass theorem \ref{thm:Stone_Weierstrass}).



\subsection{Independent Events}

\defnr{Independence of events}{Events_independence}{
\begin{enumerate}
    \item Two events $A$ and $B$ are said to be \emph{independent} if
    \[ \PP(A\cap B)=\PP(A)\,\PP(B). \]

    \item  Let $A_1,\dots,A_n$ be events. We say that the events $A_1,\dots,A_n$ are \emph{independent} if, for every nonempty subset $\{i_1,\dots,i_p\}\subset\{1,\dots,n\}$, we have
    \[ \PP(A_{i_1}\cap\cdots\cap A_{i_p}) = \PP(A_{i_1})\cdots\PP(A_{i_p}). \]
\end{enumerate}
}
\textbf{remark:} In particular, for $n\ge3$, pairwise independence does not imply independence.

\proppr{Multi_event_independence}{
The events $A_1,\dots,A_n$ are independent if and only if
\[ \PP(B_1\cap\cdots\cap B_n) = \PP(B_1)\cdots\PP(B_n) \]
for every choice of events $B_i\in\sigma(A_i) =\{\varnothing,A_i,A_i^c,\Omega\}$, $i=1,\dots,n$.
}{
    Left exercise for the reader.
}


\subsection{Independence for \sfields and Random Variables}
We say that $\G$ is a \emph{sub-\sfield} of $\A$ if
$\G$ is a \sfield and $\G\subset\A$.
Roughly speaking, a sub-\sfield represents partial information in the
probability space, namely the information obtained by knowing which events in
$\G$ have occurred. For instance, if $X$ is a random variable, the
\sfield $\sigma(X)$ corresponds to the information provided by the value
of $X$.

This suggests that the most general notion of independence is that of
independent sub-\sfields: informally, two sub-\sfields
$\G$ and $\G'$ are independent if knowing which events of
$\G$ have occurred gives no information about the occurrence of events
in $\G'$, and conversely. We make this precise in the following
definition.

\defnr{Independence of [finite] \sfields}{sigmafield_independence}{
\begin{enumerate}
\item Let $\G_1,\dots,\G_n$ be sub-\sfields of $\A$. We say that $\G_1,\dots,\G_n$ are \emph{independent} if, for every choice of events
\[ A_1\in\G_1,\ A_2\in\G_2,\ \dots,\ A_n\in\G_n, \]
we have
\[ \PP(A_1\cap A_2\cap\cdots\cap A_n) = \PP(A_1)\,\PP(A_2)\cdots\PP(A_n). \]
\item Let $X_1,\dots,X_n$ be random variables taking values in the measurable spaces $(E_1,\E_1),\dots,(E_n,\E_n)$, respectively. We say that the random variables $X_1,\dots,X_n$ are \emph{independent} if the \sfields $\sigma(X_1),\dots,\sigma(X_n)$ are independent.

\item Equivalently, $X_1,\dots,X_n$ are independent if and only if, for every choice of sets $F_1\in\E_1,\dots,F_n\in\E_n$,
\[ \PP\bigl(X_1\in F_1,\dots,X_n\in F_n\bigr) = \PP(X_1\in F_1)\cdots\PP(X_n\in F_n). \]
\end{enumerate}
}
Firstly, we introduce a suffisant condition for the independence of \sfields. 
\thmpr{Independence Extension Theorem}{extension_independence}{
Let $\G_1,\dots,\G_n$ be sub-\sfields of $\A$. For each $i\in\{1,\dots,n\}$, let $\C_i\subset\G_i$ be a class of sets such that:
\begin{itemize}
\item $\C_i$ is closed under finite intersections,
\item $\Omega\in\C_i$,
\item $\sigma(\C_i)=\G_i$.
\end{itemize}
Assume that
\[ \forall C_1\in\C_1,\dots,C_n\in\C_n, \qquad \PP(C_1\cap\cdots\cap C_n) = \prod_{i=1}^n \PP(C_i). \]
Then the \sfields $\G_1,\dots,\G_n$ are independent.
}{
We first fix $C_2\in\C_2,\dots,C_n\in\C_n$ and define
\[ \M_1 = \Bigl\{ B_1\in\G_1 : \PP(B_1\cap C_2\cap\cdots\cap C_n) = \PP(B_1)\prod_{i=2}^n \PP(C_i) \Bigr\}. \]
By assumption, $\C_1\subset\M_1$. It is straightforward to check that $\M_1$ is a monotone class. By the monotone class theorem,
\[ \M_1 \supset \sigma(\C_1)=\G_1. \]
Hence, for all $B_1\in\G_1$ and all $C_2\in\C_2,\dots,C_n\in\C_n$,
\[ \PP(B_1\cap C_2\cap\cdots\cap C_n) = \PP(B_1)\prod_{i=2}^n \PP(C_i). \]

We now fix $B_1\in\G_1$, $C_3\in\C_3,\dots,C_n\in\C_n$ and define
\[ \M_2 = \Bigl\{ B_2\in\G_2 : \PP(B_1\cap B_2\cap C_3\cap\cdots\cap C_n) = \PP(B_1)\PP(B_2)\prod_{i=3}^n \PP(C_i) \Bigr\}. \]
Again, $\C_2\subset\M_2$, and $\M_2$ is a monotone class, so $\M_2=\G_2$.

Proceeding inductively, we obtain that, for all $B_1\in\G_1,\dots,B_n\in\G_n$,
\[ \PP(B_1\cap\cdots\cap B_n) = \prod_{i=1}^n \PP(B_i). \]
This is precisely the definition of independence of the \sfields $\G_1,\dots,\G_n$.
}

The next theorem provides a fundamental characterization of the independence of $n$ random variables.

\thmpr{Independence Characterization}{independence_characterization}{
Let $X_1,\dots,X_n$ be random variables taking values in $\RR$, the following statements are
equivalent:
\begin{enumerate}
\item
The random variables $X_1,\dots,X_n$ are independent.

\item  the law of the random vector $(X_1,\dots,X_n)$ is the product of the respective laws of $X_1,\dots,X_n$, that is,
\[ \PP_{(X_1,\dots,X_n)} = \PP_{X_1}\otimes\cdots\otimes\PP_{X_n}.\]

\item For every $a_1,\dots,a_n\in\RR$,
\[ \PP(X_1\le a_1,\dots,X_n\le a_n) = \prod_{i=1}^n \PP(X_i\le a_i). \]

\item
For every choice of continuous functions $f_1,\dots,f_n:\RR\to\RR_+$ with compact support,
\[ \EE\!\left[\prod_{i=1}^n f_i(X_i)\right] = \prod_{i=1}^n \EE[f_i(X_i)]. \]

\item The characteristic function of the random vector $X=(X_1,\dots,X_n)$ factorizes as
\[ \Phi_X(\xi_1,\dots,\xi_n) = \prod_{i=1}^n \Phi_{X_i}(\xi_i), \qquad (\xi_1,\dots,\xi_n)\in\RR^n. \]
\end{enumerate}
}{
\begin{itemize}
    \item $(1)\Longleftrightarrow(2)$ : \\ 
Let $F_i \in \E_i$, for every $i \in \{1,\ldots,n\}$. On one hand,
\[ \PP_{(X_1,\ldots,X_n)}(F_1 \times \cdots \times F_n) = \PP\big( (X_1 \in F_1) \cap \cdots \cap (X_n \in F_n) \big). \]
On the other hand,
\[ \PP_{X_1} \otimes \cdots \otimes \PP_{X_n}(F_1 \times \cdots \times F_n) = \prod_{i=1}^n \PP_{X_i}(F_i) = \prod_{i=1}^n \PP(X_i \in F_i). \]
Comparing with the definition of independence, we see that $X_1,\ldots,X_n$ are independent if and only if the two probability measures $\PP_{(X_1,\ldots,X_n)}$ and $\PP_{X_1} \otimes \cdots \otimes \PP_{X_n}$ take the same values on boxes $F_1 \times \cdots \times F_n$.
This is equivalent to saying that : 
\[ \PP_{(X_1,\ldots,X_n)} = \PP_{X_1} \otimes \cdots \otimes \PP_{X_n}, \]
since we know that a probability measure on a product space is characterized by its values on boxes (theorem \ref{thm:product_measure_existence}) we conclude the result.
    \item $(2)\Longleftrightarrow(3)$ : \\ 
Just by taking $F_i=]-\infty,a_i]$ we can prove the $(2)\Longleftarrow (3)$, conversely,
Just by taking $C_i = \bigl\{ ]-\infty, a] \;\big|\; a \in \RR \bigr\}$, and apply the theorem \ref{thm:suffisant_condition_sigmafield_independence} we obtain the independence $(1)$ and therefore $(2)$.
    \item $(1)\Longleftrightarrow(4)$ : \\ 
We have $(4) \Longrightarrow (3) \Longrightarrow (1)$ by taking $f_i(x) = \1_{x\leq a_i}$.
The other implication is imediate from the next proposition. 

    \item $(2)\Longleftrightarrow(5)$ : \\  To show this equivalence, we note that the Fourier transform of the product measure $\PP_{X_1} \otimes \cdots \otimes \PP_{X_n}$ is
\begin{align*}
(\xi_1,\ldots,\xi_n)
&\longmapsto
\int \exp\!\left( i \sum_{j=1}^n \xi_j x_j \right)
\, \PP_{X_1}(dx_1) \cdots \PP_{X_n}(dx_n) \\
&= \prod_{j=1}^n \int e^{i \xi_j x_j}\, \PP_{X_j}(dx_j) \\
&= \prod_{j=1}^n \Phi_{X_j}(\xi_j).
\end{align*}

using the Fubini theorem in the first equality. Hence the property (5) is equivalent to the fact that the Fourier transform of $\PP_{(X_1,\ldots,X_n)}$ coincides with the Fourier transform of $\PP_{X_1} \otimes \cdots \otimes \PP_{X_n}$. By pervious property, this holds if and only if
\[ \PP_{(X_1,\ldots,X_n)} = \PP_{X_1} \otimes \cdots \otimes \PP_{X_n}, \]

\end{itemize}
}

From this theorem follow this interesting results :
\proppr{Random_variables_Independence_proprieties}{
Let $X_1,\dots,X_n$ be independent real-valued random variables.
\begin{enumerate}
\item
Let $f_1,\dots,f_n$ be non negative measurable functions we have : 
\[ \EE\!\left[\prod_{i=1}^n f_i(X_i)\right] = \prod_{i=1}^n \EE\!\left[f_i(X_i)\right], \]
\item
Let $f_1,\dots,f_n$ be measurable functions such that
\[ \EE[|f_i(X_i)|]<\infty \quad\text{for every } i\in\{1,\dots,n\}. \]
Then
\[ \EE\!\left[\prod_{i=1}^n f_i(X_i)\right] = \prod_{i=1}^n \EE[f_i(X_i)], \]
and the expectation on the left-hand side is well defined. The same statement holds for complex-valued functions $f_i$ under the same integrability assumption.

\item As a special case of (1), if $X_1,\dots,X_n$ are independent random variables in $\L^1$, then $X_1\cdots X_n\in \L^1$ and
\[ \EE[X_1\cdots X_n] = \prod_{i=1}^n \EE[X_i]. \]
In general, the product of random variables in $\L^1$ need not belong to $\L^1$ without independence.

\item If $X_1$ and $X_2$ are independent real-valued random variables in $L^2$, then
\[ \operatorname{cov}(X_1,X_2)=0. \]

\item Suppose that, for every $i\in\{1,\dots,n\}$, the random variable $X_i$ admits a density $p_i$ with respect to Lebesgue measure and that $X_1,\dots,X_n$ are independent. Then the random vector $(X_1,\dots,X_n)$ admits a density $p$ given by
\[ p(x_1,\dots,x_n)=\prod_{i=1}^n p_i(x_i). \]

\item Conversely, suppose that $(X_1,\dots,X_n)$ admits a density $p$ of the form
\[ p(x_1,\dots,x_n)=\prod_{i=1}^n q_i(x_i), \]
where $q_1,\dots,q_n$ are nonnegative measurable functions on $\RR$. Then $X_1,\dots,X_n$ are independent, and for every $i\in\{1,\dots,n\}$, $X_i$ admits a density $p_i$ of the form $p_i=C_i q_i$, where $C_i>0$ is a normalization constant.
\end{enumerate}
}{
\begin{enumerate}
    \item  Using the Fubini-Tonelli theorem \ref{thm:Fubini_Tonelli} for the case of two independent variables : 
    \begin{align*}
    \EE\bigl(f_1(X_1)f_2(X_2)\bigr)
            &= \int_{\RR^2} f_1(x_1)f_2(x_2), d\PP_{(X_1,X_2)}(x_1,x_2) \\
            &= \int_{\RR^2} f_1(x_1)f_2(x_2), d(\PP_{X_1}\otimes \PP_{X_2})(x_1,x_2) \\
            &= \int_{\RR} \int_{\RR} f_1(x_1)f_2(x_2), d\PP_{X_2}(x_2), d\PP_{X_1}(x_1) \\
            &= \int_{\RR} f_1(x_1), d\PP_{X_1}(x_1)
            \int_{\RR} f_2(x_2), d\PP_{X_2}(x_2) \\
            &= \EE\bigl(f_1(X_1)\bigr)\times \EE\bigl(f_2(X_2)\bigr).
    \end{align*}
    The general case can be achived by a simple reccurence. 
    \item Similar to (1), but using the Fubini-Lebesgue theorem \ref{thm:Fubini_Lebesgue} and ensuring the $f_i(X_i)\in \L^1$ which is the condition :
    \[ \EE[|f_i(X_i)|]<\infty \quad\text{for every } i\in\{1,\dots,n\}. \]
    \item Clear conclusion from (1). 
    \item This follows from the preceding remarks since :
    \[ \operatorname{Cov}(X_1,X_2) = \EE(X_1X_2) - \EE(X_1)\EE(X_2) = 0 \] 
    \item This is an immediate consequence of the caracterisation of independece $\PP_{(X_1,\ldots,X_n)} = \PP_{X_1} \otimes \cdots \otimes \PP_{X_n}$. Indeed, by the Fubini theorem, under the assumption
    \[ \PP_{X_i}(dx_i) = p_i(x_i)\,dx_i \quad \text{for } i \in \{1,\ldots,n\}, \]
    we have
    \[ \PP_{X_1} \otimes \cdots \otimes \PP_{X_n}(dx_1\cdots dx_n) = \left( \prod_{i=1}^n p_i(x_i) \right) dx_1\cdots dx_n. \]

    \item We first observe, again thanks to the Fubini theorem, that
    \[ \prod_{i=1}^n \left( \int q_i(x)\,dx \right) = \int_{\RR^n} p(x_1,\ldots,x_n)\,dx_1\cdots dx_n = 1. \]
    In particular, the quantity
    \[K_i = \int q_i(x)\,dx\]
    is positive and finite for every \( i \in \{1,\ldots,n\} \).
    Then, by Properties of density function \ref{defn:density_function}, the density of \( X_i \) is
    \begin{align*}
    p_i(x_i)
        &= \int_{\RR^{n-1}} p(x_1,\ldots,x_n)\, dx_1\cdots dx_{i-1}\,dx_{i+1}\cdots dx_n \\
        &= \left( \prod_{j\neq i} K_j \right) q_i(x_i) = \frac{1}{K_i}\, q_i(x_i).
    \end{align*}
    Thus \( p_i = C_i q_i \) with \( C_i = 1/K_i \).
    This also allows us to rewrite the density of \( (X_1,\ldots,X_n) \) as
    \[ p(x_1,\ldots,x_n) = \prod_{i=1}^n q_i(x_i) = \prod_{i=1}^n p_i(x_i), \]
    and therefore
    \[ \PP_{(X_1,\ldots,X_n)} = \PP_{X_1} \otimes \cdots \otimes \PP_{X_n}, \]
    which shows that \( X_1,\ldots,X_n \) are independent.
\end{enumerate}
}


\defnr{Independence of [infinite] \sfields}{infinite_sigmafield_independence}{
Let $(\G_i)_{i\in I}$ be an arbitrary collection of sub-\sfields of
$\A$. We say that the \sfields $\G_i$, $i\in I$, are
\emph{independent} if, for every finite subset
$\{i_1,\dots,i_p\}\subset I$, the \sfields
$\G_{i_1},\dots,\G_{i_p}$ are independent.

If $(X_i)_{i\in I}$ is an arbitrary collection of random variables, we say that
the random variables $X_i$, $i\in I$, are \emph{independent} if the
\sfields $\sigma(X_i)$, $i\in I$, are independent.

Similarly, for an arbitrary collection $(A_i)_{i\in I}$ of events, we say that
the events $A_i$, $i\in I$, are independent if the \sfields
$\sigma(A_i)$, $i\in I$, are independent.
}


\proppr{sigmafield_independence_separation}{
Let $(X_n)_{n\in\NN}$ be a sequence of independent random variables. Then,
for every $p\in\NN$, the \sfields
\[ 
\G_1=\sigma(X_1,\dots,X_p),
\qquad
\G_2=\sigma(X_{p+1},X_{p+2},\dots)
\]
are independent.
}{
We apply the suffisant condition of \sfields independence given in \ref{thm:suffisant_condition_sigmafield_independence}.
Set
\[ \C_1=\sigma(X_1,\dots,X_p)=\G_1 \]
and
\[ \C_2=\bigcup_{k=p+1}^{\infty} \sigma(X_{p+1},\dots,X_k) \subset \G_2. \]
The class $\C_2$ is closed under finite intersections, contains $\Omega$, and satisfies $\sigma(\C_2)=\G_2$.

By independence of the random variables $(X_n)_{n\in\NN}$, and by the grouping-by-blocks principle, we have for every $C_1\in\C_1$ and $C_2\in\C_2$,
\[ \PP(C_1\cap C_2)=\PP(C_1)\PP(C_2). \]
Therefore, the assumptions of theorem \ref{thm:suffisant_condition_sigmafield_independence} are satisfied, and it follows that the \sfields $\G_1$ and $\G_2$ are independent.
}

\subsection{Sums of Independent Random Variables}
Sums of independent random variables play a central role in probability theory. Their distributional and analytic properties are conveniently described using the notion of convolution of probability measures.
Let $\PP$ and $\QQ$ be two probability measures on $\RR^d$. The \emph{convolution} $\PP*\QQ$ is defined as the pushforward of the product measure $\PP\otimes\QQ$ under the mapping $(x,y)\mapsto x+y$. Equivalently, for every nonnegative measurable function $\varphi$ on $\RR^d$,
\[ \int_{\RR^d} \varphi(z)\,(\PP*\QQ)(dz) = \int_{\RR^d}\!\!\int_{\RR^d} \varphi(x+y)\,d\PP(x)\,d\QQ(y). \]

If $\PP$ and $\QQ$ admit densities $f$ and $g$ with respect to Lebesgue measure, then $\PP*\QQ$ admits the density $f*g$ given by
\[ (f*g)(z)=\int_{\RR^d} f(x)\,g(z-x)\,dx. \]

\proppr{convolution_independence}{
Let $X$ and $Y$ be two independent random variables with values in $\RR^d$.
\begin{enumerate}
\item The law of $X+Y$ is the convolution of the laws of $X$ and $Y$:
\[ \PP_{X+Y}=\PP_X*\PP_Y. \]
In particular, if $X$ and $Y$ admit densities $p_X$ and $p_Y$, then $X+Y$ admits the density $p_X*p_Y$.

\item The characteristic function of $X+Y$ satisfies
\[ \Phi_{X+Y}(\xi)=\Phi_X(\xi)\,\Phi_Y(\xi), \qquad \xi\in\RR^d. \]

\item If $\EE[\|X\|^2]<\infty$ and $\EE[\|Y\|^2]<\infty$, then the covariance matrices satisfy
\[ K_{X+Y}=K_X+K_Y. \]
In particular, if $d=1$ and $X,Y\in L^2$, then
\[ \operatorname{var}(X+Y)=\operatorname{var}(X)+\operatorname{var}(Y). \]
\end{enumerate}
}{
\begin{enumerate}
\item We know that $\PP_{(X,Y)} = \PP_X \otimes \PP_Y$, and thus, for every nonnegative measurable function $\varphi$ on $\RR^d$,
\[ \EE[\varphi(X+Y)] = \int \! \int \varphi(x+y)\,\PP_X(dx)\,\PP_Y(dy) = \int \varphi(z)\,(\PP_X * \PP_Y)(dz), \]
by the definition of $\PP_X * \PP_Y$. The case with density follows from the remarks before the proposition.

\item This follows from the equality
\[\PP_{X+Y} = \PP_X * \PP_Y\]
and the last observation before the proposition.

\item If $X = (X_1,\ldots,X_d)$ and $Y = (Y_1,\ldots,Y_d)$, proposition \ref{prop:Random_variables_Independence_proprieties} implies that
\[ \operatorname{cov}(X_i,Y_j) = 0 \quad \text{for every } i,j \in \{1,\ldots,d\}. \]
Consequently, using bilinearity,
\[ \operatorname{cov}(X_i+Y_i,\,X_j+Y_j) = \operatorname{cov}(X_i,X_j) + \operatorname{cov}(Y_i,Y_j), \]
which gives $ K_{X+Y} = K_X + K_Y. $
\end{enumerate}
}


\subsection{The Borel--Cantelli Lemma}
Let $(A_n)_{n\in\NN}$ be a sequence of events. Recall that
\[ \limsup_{\goto{n}} A_n = \bigcap_{n=1}^{\infty}\bigcup_{k=n}^{\infty} A_k. \]
The event $\limsup_{\goto{n}} A_n$ consists of those $\omega\in\Omega$ that belong to infinitely many of the events $A_n$. Note that It's easy to check that :
\[ \PP\bigl(\limsup_{\goto{n}} A_n\bigr) \ge \limsup_{\goto{n}} \PP(A_n). \]

\thmpr{Borel--Cantelli Lemma}{Borel_Cantelli}{
Let $(A_n)_{n\in\NN}$ be a sequence of events.
\begin{enumerate}
\item If $\sum_{n\in\NN} \PP(A_n)<\infty,$ then
\[ \PP\bigl(\limsup_{\goto{n}} A_n\bigr)=0. \]
Equivalently,
\[ \bigl\{n\in\NN:\ \omega\in A_n\bigr\} \quad\text{is finite almost surely.} \]

\item
If $\sum_{n\in\NN} \PP(A_n)=\infty$ and the events $(A_n)_{n\in\NN}$ are independent, then
\[ \PP\bigl(\limsup_{\goto{n}} A_n\bigr)=1. \]
Equivalently,
\[ \bigl\{n\in\NN:\ \omega\in A_n\bigr\} \quad\text{is infinite almost surely.} \]
\end{enumerate}
}{
\begin{enumerate}
\item If $\sum_{n\in\NN} \PP(A_n)<\infty$, then
\[ \EE\!\left[\sum_{n\in\NN} \1_{A_n}\right] = \sum_{n\in\NN} \PP(A_n) < \infty. \]
It follows that $ \sum_{n\in\NN} \1_{A_n}<\infty \quad\text{almost surely}, $
which means that $\omega$ belongs to only finitely many of the events $A_n$. Hence,
\[ \PP\bigl(\limsup_{\goto{n}} A_n\bigr)=0. \]

\item
Assume now that $\sum_{n\in\NN} \PP(A_n)=\infty$ and that the events $A_n$ are independent. Fix $n_0\in\NN$. For $n\ge n_0$, we have
\[ \PP\!\left(\bigcap_{k=n_0}^{n} A_k^c\right) = \prod_{k=n_0}^{n} \PP(A_k^c) = \prod_{k=n_0}^{n} \bigl(1-\PP(A_k)\bigr). \]
Since the series $\sum_{k\in\NN} \PP(A_k)$ diverges, the right-hand side converges to $0$ as $\goto{n}$. Consequently, $ \PP\!\left(\bigcap_{k=n_0}^{\infty} A_k^c\right)=0$. As this holds for every $n_0\in\NN$, we obtain
\[ \PP\!\left(\bigcup_{n_0=1}^{\infty} \bigcap_{k=n_0}^{\infty} A_k^c\right)=0. \]
Taking complements yields
\[ \PP\!\left(\bigcap_{n_0=1}^{\infty} \bigcup_{k=n_0}^{\infty} A_k\right)=1, \]
that is, $ \PP\bigl(\limsup_{\goto{n}} A_n\bigr)=1. $
\end{enumerate}
}

The Borel--Cantelli Lemma is a powerful tool for proving (or disproving) that a certain property holds eventually almost surely, meaning that from an index onward the property is true almost surely. It achieves this by reducing the problem to studying whether a series of probabilities converges or diverges. The standard application of the lemma is summarized in Figure~\ref{fig:borel_cantelli_strategy}.

\begin{figure}[H]
\centering
\begin{tikzpicture}[
    node distance=1cm,
    >=Latex,
    every node/.style={
        draw,
        rounded corners=2pt,
        align=center,
        text width=15cm,
        minimum height=1cm,
        font=\small
    },
    every edge/.style={draw, thick, ->}
]

\node (step1) {
\textbf{Step 1. Identify the bad event.}\\[2mm]
Define $ A_n=\{\text{the desired property fails at time }n\}.$
};

\node (step2) [below=of step1] {
\textbf{Step 2. Estimate the probabilities.}\\[2mm]
Find an upper bound for $P(A_n)$
Usually use Markov, Chebyshev, Hoeffding, Chernoff, or another inequality.
};

\node (step3) [below=of step2] {
\textbf{Step 3. Check the series.}\\[2mm]
Determine whether $\sum_{n=1}^{\infty}P(A_n)$ converges or diverges.
};

\node (step4) [below=of step3] {
\textbf{Step 4. Apply Borel--Cantelli.}\\[2mm]
If $\sum P(A_n)<\infty,$ then $P(A_n\ \text{i.o.})=0.$ \\
so eventually the desired property always holds.
};

\draw (step1)--(step2);
\draw (step2)--(step3);
\draw (step3)--(step4);

\end{tikzpicture}
\caption{Standard strategy for applying the First Borel--Cantelli Lemma.}
\label{fig:borel_cantelli_strategy}
\end{figure}

A more general version of the second lemma of Borel--Cantelli is the zero-one law of kolmogrov that stated as follow :
\thmr{Kolmogorov's Zero--One Law}{zero_one_law}{
Let $(\Omega,\F,\PP)$ be a probability space, and let
$(\F_n)_{n\ge 1}$ be a sequence of independent sub-$\sigma$-algebras
of $\F$.
For each $n\ge 1$, define the tail $\sigma$-algebra from time $n$ onward by
\[ \B_n := \sigma\!\left(\bigcup_{k=n}^{\infty}\F_k\right). \]

The \emph{tail} (or \emph{terminal}) $\sigma$-algebra associated with
$(\F_n)_{n\ge 1}$ is defined by
\[ \B_\infty := \bigcap_{n=1}^{\infty}\B_n. \]

Then the tail $\sigma$-algebra is $\PP$-trivial. That is,
\[ \forall A\in\B_\infty, \qquad \PP(A)\in\{0,1\}. \]
}
Let $(X_n)_{n\ge 1}$ be a sequence of independent random variables on
$(\Omega,\F,\PP)$. We can state the zero-one law as follow : consider the associated tail $\sigma$-algebra is \[ \B_\infty = \bigcap_{n=1}^{\infty} \sigma(X_n,X_{n+1},\ldots). \]
For every event $A\in\B_\infty$,
\[ \PP(A)\in\{0,1\}. \]

\section{Conditioning}

\subsection{Definition of conditional expectaiton} 
\subsubsection{Discret case}
As in the previous chapters, we consider a probability space $(\probspace)$. If $B \in \A$ is an event of positive probability,
we can define a new probability measure on $(\measspace)$ by setting, for every $A \in \A$, \[ \PP(A \mid B) := \frac{\PP(A \cap B)}{\PP(B)}. \]

This probability measure $A \mapsto \PP(A \mid B)$ is called \textbf{the conditional probability given $B$}. Similarly, for every nonnegative random variable $X$, or for $X \in L^{1}(\probspace)$,
the \textbf{conditional expectation of $X$ given $B$} is defined by
\[ \EE[X \mid B] := \frac{\EE[X \1_B]}{\PP(B)}. \]

One verifies that $\EE[X \mid B]$ is also the expected value of $X$ under $\PP(\,\cdot \mid B)$.
We can interpret $\EE[X \mid B]$ as the mean value of $X$ knowing that the event $B$ occurs.

Let us now define the conditional expectation knowing a discrete random variable.
We consider a random variable $Y$ taking values in a countable space $E$ equipped as usual
with the $\sigma$-field of all its subsets.
Let
\[ E' = \{ y \in E : \PP(Y = y) > 0 \}. \]
If $X \in L^{1}(\Omega, \A, \PP)$, we may consider, for every $y \in E'$,
\[ \EE[X \mid Y = y] = \frac{\EE[X \1_{\{Y=y\}}]}{\PP(Y = y)}, \]
as a special case of the conditional expectation $\EE[X \mid B]$ when $\PP(B) > 0$.

The conditional expectation of $X$ knowing $Y$ is the real random variable
\[\EE[X \mid Y] = \varphi(Y),\]
where the function $\varphi : E \to \RR$ is given by
\[\varphi(y) =
\begin{cases}
\EE[X \mid Y = y], & \text{if } y \in E',\\
0, & \text{if } y \in E \setminus E'.
\end{cases} \]

\subsubsection{General case}
Now we turn to the definition of conditional expectation with respect to a sub-\sfield. For that purpose we need the next theorem that provides the existance and uniqueness of conditional expectation.
\thmpr{Existance of Conditional Expectation}{Conditional_expectation_existence}{
Let $\B$ be a sub-$\sigma$-field of $\A$, and let $X \in L^1(\Omega,\A,\PP)$. There exists a unique element of $L^1(\Omega,\B,\PP)$, which is denoted by $\EE[X \mid \B]$, such that
\begin{equation} \label{eq:conditional_expectation_characteristic_property1}\forall B \in \B, \quad \EE[X \1_B] = \EE[\EE[X \mid \B] \1_B]. \tag{i} 
\end{equation}

We have more generally, for every bounded $\B$-measurable real random variable $Z$,
\begin{equation} \label{eq:conditional_expectation_characteristic_property2}
\EE[XZ] = \EE[\EE[X \mid \B] Z]. \tag{ii}
\end{equation}
If $X \ge 0$, we have $\EE[X \mid \B] \ge 0$ a.s.
}{
Let us start by proving uniqueness. Let $X'$ and $X''$ be two random variables in $L^1(\Omega,\B,\PP)$ such that
\[ \forall B \in \B, \quad \EE[X' \1_B] = \EE[X \1_B] = \EE[X'' \1_B]. \]
Taking $B = \{X' > X''\}$ (which is in $\B$ since both $X'$ and $X''$ are $\B$-measurable), we get
\[ \EE[(X' - X'') \1_{\{X' > X''\}}] = 0, \]
which implies that $X' \le X''$ a.s., and we have similarly $X' \ge X''$ a.s. Thus $X' = X''$ a.s., which means that $X'$ and $X''$ are equal as elements of $L^1(\Omega,\B,\PP)$.

Let us now turn to existence. We first assume that $X \ge 0$ and $\EE(X)\neq 0$, then we let $\QQ$ be the finite measure on $(\Omega,\B)$ defined by
\[ \forall B \in \B, \quad \QQ(B) := \frac{1}{\EE(X)}\EE[X \1_B]. \]

Let us emphasize that we define $\QQ(B)$ only for $B \in \B$. We may also view $\PP$ as a probability measure on $(\Omega,\B)$, by restricting the mapping $B \mapsto \PP(B)$ to $B \in \B$, and it is immediate that $\QQ \ll \PP$. The Radon--Nikodym theorem \ref{thm:RadonNikodym} applied to the probability measures $\PP$ and $\QQ$ on the measurable space $(\Omega,\B)$ yields the existence of a nonnegative $\B$-measurable random variable $\tilde X$ such that
\[ \forall B \in \B, \quad \EE[X \1_B] = \QQ(B) = \EE[\tilde X \1_B]. \]
Taking $B = \Omega$, we have $\EE[\tilde X] = \EE[X] < \infty$, and thus $\tilde X \in L^1(\Omega,\B,\PP)$. The random variable $\EE[X \mid \B] = \tilde X$ satisfies (1). (for the case of $\EE(X)=0$ we have $X=0$ a.s. and we can take $\tilde{X}=0$)

When $X$ is of arbitrary sign, we just have to take
\[ \EE[X \mid \B] = \EE[X^+ \mid \B] - \EE[X^- \mid \B], \]
and it is clear that (1) also holds in that case.

Finally, to see that (1) implies (2), we rely on the usual measure-theoretic arguments. (2) follows from (1) when $Z$ is a simple random variable (taking only finitely many values), and in the general case lemma \ref{lem:approximation_simple_random_variables} allows us to write $Z$ as the pointwise limit of a sequence $(Z_n)_{n \in \NN}$ of simple $\B$-measurable random variables that are uniformly bounded by the same constant $K$ (such that $|Z| \le K$) and the dominated convergence theorem yields the desired result.
}

The crucial point is the fact that $\EE[X \mid \B]$ is $\B$-measurable. Either of properties (1) and (2) characterizes the conditional expectation $\EE[X \mid \B]$ among random variables of $L^1(\Omega,\B,\PP)$. In what follows, we will refer to (1) or (2) as the \emph{characteristic property} of $\EE[X \mid \B]$.

In the special case where $\B = \sigma(Y)$ is generated by a random variable $Y$, we will write indifferently
\[ \EE[X \mid \B] = \EE[X \mid \sigma(Y)] = \EE[X \mid Y]. \]


We now turn to the definition of $\EE[X \mid \B]$ for a nonnegative random variable $X$ which allows $X$ to take the value $+\infty$.

\thmpr{Existence of Conditional Expectation for nonnegative random variables}{Conditional_expectation_existence_nonnegative}{
Let $X$ be a random variable with values in $[0,\infty]$. There exists a $\B$-measurable random variable with values in $[0,\infty]$, which is denoted by $\EE[X \mid \B]$ and is such that, for every nonnegative $\B$-measurable random variable $Z$,
\[ \EE[XZ] = \EE[\EE[X \mid \B] Z]. \tag{iii} \]
Furthermore $\EE[X \mid \B]$ is unique up to a $\B$-measurable set of probability zero.
}{
We define $\EE[X \mid \B]$ by setting
\[ \EE[X \mid \B] := \lim_{n \uparrow \infty} \EE[X \wedge n \mid \B] \quad \text{a.s.} \]
This definition makes sense because, for every $n \in \NN$, $X \wedge n$ is bounded hence integrable and thus $\EE[X \wedge n \mid \B]$ is well defined by the previous section. Furthermore, the fact that the sequence $(\EE[X \wedge n \mid \B])_{n \in \NN}$ is (a.s.) increasing follows from the fact that the conditional expectation of a.s. non-negative random variable is also an a.s. non-negative random variable. Then, if $Z$ is nonnegative and $\B$-measurable, the monotone convergence theorem implies that
\[ \EE[\EE[X \mid \B] Z] = \lim_{n \to \infty} \EE[\EE[X \wedge n \mid \B] Z \1_{\{Z < \infty\}}] = \lim_{n \to \infty} \EE[(X \wedge n) Z \1_{\{Z < \infty\}}] = \EE[XZ]. \]

It remains to establish uniqueness. Let $X'$ and $X''$ be two $\B$-measurable random variables with values in $[0,\infty]$, such that
\[ \EE[X'Z] = \EE[X''Z] \]
for every nonnegative $\B$-measurable random variable $Z$. Let us fix two nonnegative rationals $a < b$, and take
\[ Z = \1_{\{X' \le a < b \le X''\}}. \]

It follows that
\[ a \, \PP(X' \le a < b \le X'') \le \EE[\1_{\{X' \le a < b \le X''\}} X'] = \EE[\1_{\{X' \le a < b \le X''\}} X''] \ge b \, \PP(X' \le a < b \le X''), \]
which is only possible if $\PP(X' \le a < b \le X'') = 0$. Hence,
\[ \PP\!\left( \bigcup_{\substack{a,b \in \QQ_+ \\ a < b}} \{ X' \le a < b \le X'' \} \right) = 0, \]
which implies $X' \ge X''$ a.s., and interchanging the roles of $X'$ and $X''$ also gives $X'' \ge X'$ a.s.
}

\textbf{Remark:}
The uniqueness part of the theorem means that we should view $\EE[X \mid \B]$ as an equivalence class of $\B$-measurable random variables that are equal a.s. But of course it will be more convenient to speak about "the random variable" $\EE[X \mid \B]$.

In the case where $X$ is also integrable, the definition of the preceding theorem is consistent with that of Theorem \ref{thm:Conditional_expectation_existence}. Indeed (3) with $Z = 1$ shows that $\EE[\EE[X \mid \B]] = \EE[X]$, and in particular $\EE[X \mid \B] < \infty$ a.s.\ so that we can take a representative with values in $[0,\infty)$ and $\EE[X \mid \B] \in L^1$. We then just have to note that (1) (in the theorem we are citing) is the special case of (3) where $Z = \1_B$, $B \in \B$.

Similarly as in the case of integrable variables, we will refer to (3) (or to its variant when $Z = \1_B$, $B \in \B$) as the characteristic property of $\EE[X \mid \B]$.

\subsection{Properties of conditional Expectations}
We now list the properties of conditional expectaitons : 
\proppr{conditional_expectation_proprieties}{
\begin{enumerate}
\item  For a $X,Y \in L^1(\Omega,\A,\PP)$)  and $\alpha,\beta\in\RR$ we have : 
\[ \EE(\alpha X + \beta Y \mid \B ) = \alpha \EE(X\mid \B) + \beta\EE(Y\mid \B) \] 
\item  For a $X,Y$ non negative random variables and $\alpha,\beta\in\RR^+$ we have : 
\[ \EE(\alpha X + \beta Y \mid \B ) = \alpha \EE(X\mid \B) + \beta\EE(Y\mid \B) \] 
\item If $X \in L^1(\Omega,\A,\PP)$) (or if $X$ is just non negative) and $X$ is $\B$-measurable, then :
\[\EE[X \mid \B] = X\]
\item If $X \in L^1(\Omega,\A,\PP)$) (or if $X$ is just non negative), then :
\[\EE[\EE[X \mid \B]] = \EE[X]\]
\item If $X, X' \in L^1(\Omega,\A,\PP)$) and $X \ge X'$, then : 
\[ \EE[X \mid \B] \ge \EE[X' \mid \B] \quad \text{a.s} \]
\item If $X \in L^1(\Omega,\A,\PP)$), then :
\[ |\EE[X \mid \B]| \le \EE[|X| \mid \B] \quad \text{a.s.}\] 
and, consequently,
\[ \EE[|\EE[X \mid \B]|] \le \EE[|X|]. \]
Therefore the mapping $X \mapsto \EE[X \mid \B]$ is a contraction of $L^1(\Omega,\A,\PP)$).
\end{enumerate}
}{
    The proofs for these properties are left exercise for the reader.
}
Now the properties allowing intereaction between limits and conditional expectations : 
\proppr{conditional_expectation_limits_properties}{
\begin{enumerate}
\item If $(X_n)_{n \in \NN}$ is an increasing sequence of nonnegative random variables, and $X = \lim \uparrow X_n$, then
\[ \EE[X \mid \B] = \lim_{n \to \infty} \uparrow \, \EE[X_n \mid \B], \quad \text{a.s.} \]

As consequence : if $(Y_n)_{n \in \NN}$ is a sequence of nonnegative random variables, we have
\[ \EE\!\left[ \sum_{n \in \NN} Y_n \,\middle|\, \B \right] = \sum_{n \in \NN} \EE[Y_n \mid \B]. \]

\item If $(X_n)_{n \in \NN}$ is any sequence of nonnegative random variables, then
\[ \EE[\liminf X_n \mid \B] \le \liminf \EE[X_n \mid \B], \quad \text{a.s.} \]

\item Let $(X_n)_{n \in \NN}$ be a sequence of integrable random variables that converges a.s.\ to $X$. Assume that there exists a nonnegative random variable $Z$ such that $|X_n| \le Z$ a.s.\ for every $n \in \NN$, and $\EE[Z] < \infty$. Then,
\[ \EE[X \mid \B] = \lim_{n \to \infty} \EE[X_n \mid \B], \quad \text{a.s.}  \]
The convergence is also in $\L^1$ that is : 
\[ \EE\left[ \big| \EE[X \mid \B] - \EE[X_n \mid \B] \big| \right] \longrightarrow 0 \quad \textit{as } n\rightarrow +\infty  \]
\end{enumerate}
}{
\begin{enumerate}
\item It follows from last property that we have $\EE[X \mid \B] \ge \EE[Y \mid \B]$ if $X \ge Y \ge 0$. Under the assumptions of (1), we can therefore set
\[ X' = \lim \uparrow \EE[X_n \mid \B], \]
which is a $\B$-measurable random variable with values in $[0,\infty)$. Then, for every nonnegative $\B$-measurable random variable $Z$, the monotone convergence theorem gives
\[ \EE[ZX'] = \lim \uparrow \EE[Z \EE[X_n \mid \B]] = \lim \uparrow \EE[Z X_n] = \EE[ZX], \]
which implies $X' = \EE[X \mid \B]$ thanks to the characteristic property (3) from theorem \ref{thm:Conditional_expectation_existence}. The second assertion in (1) follows by applying the first one to $X_n = Y_1 + \cdots + Y_n$.

\item Using (1), we have
\[ \EE[\liminf X_n \mid \B] = \EE\!\left[ \lim_{k \uparrow \infty} \left( \inf_{n \ge k} X_n \right) \,\middle|\, \B \right] = \lim_{k \uparrow \infty} \EE\!\left[ \inf_{n \ge k} X_n \,\middle|\, \B \right]. \]

\item It suffices to apply (2) twice:
\[ \EE[Z - X \mid \B] = \EE[\liminf (Z - X_n) \mid \B] \le \EE[Z \mid \B] - \limsup \EE[X_n \mid \B], \]
\[ \EE[Z + X \mid \B] = \EE[\liminf (Z + X_n) \mid \B] \le \EE[Z \mid \B] + \liminf \EE[X_n \mid \B], \]
which leads to
\[ \EE[X \mid \B] \le \liminf \EE[X_n \mid \B] \le \limsup \EE[X_n \mid \B] \le \EE[X \mid \B], \]
giving the desired almost sure convergence. The convergence in $L^1$ is now a consequence of the dominated convergence theorem, since we have
\[ |\EE[X_n \mid \B]| \le \EE[|X_n| \mid \B] \le \EE[Z \mid \B] \]
and $\EE[\EE[Z \mid \B]] = \EE[Z] < \infty$.
\end{enumerate}
}

\thmpr{Jensen Inequality conditional Expectation}{conditional_Jensen_inequality}{
If $f : \RR \to \RR_+$ is convex, and if $X \in L^1$, then
\[ \EE[f(X) \mid \B] \ge f(\EE[X \mid \B]). \]
}{
Set
\[ E_f = \{ (a,b) \in \RR^2 : \forall x \in \RR, \ f(x) \ge ax + b \}. \]
Then,
\[ \forall x \in \RR, \quad f(x) = \sup_{(a,b) \in E_f} (ax + b) = \sup_{(a,b) \in E_f \cap \QQ^2} (ax + b). \]

We can take advantage of the fact that $\QQ^2$ is countable to discard a countable collection of sets of probability zero and to get that, a.s.,
\[ \EE[f(X) \mid \B] = \EE\!\left[ \sup_{(a,b) \in E_f \cap \QQ^2} (aX + b) \,\middle|\, \B \right] \ge \sup_{(a,b) \in E_f \cap \QQ^2} \EE[aX + b \mid \B] = f(\EE[X \mid \B]). \quad  \]
}
As a consequence of Jensen's inequality, we obtain that, for every $p \ge 1$, the mapping
$X \mapsto \EE[X \mid \B]$ maps $L^p(\Omega,\A,\PP)$ into itself, and is a contraction of $L^p(\Omega,\A,\PP)$. Indeed, for every $X \in L^p(\Omega,\A,\PP)$, we have using property (6) in Proposition \ref{prop:conditional_expectation_proprieties}  ,
\[ \EE[|\EE[X \mid \B]|^p] \le \EE[\EE[|X| \mid \B]^p] \le \EE[\EE[|X|^p \mid \B]] = \EE[|X|^p]. \]

This contraction property also holds for $p = \infty$.

\thmpr{The orthogonal projection property}{Orthogonal_projection_property_conditional_expectation}{
If $X \in L^2(\Omega,\A,\PP)$, then $\EE[X \mid \B]$ is the orthogonal projection of $X$ on $L^2(\Omega,\B,\PP)$.
}{
Jensen's inequality shows that $\EE[X \mid \B]^2 \le \EE[X^2 \mid \B]$, a.s. This implies that
\[
\EE[\EE[X \mid \B]^2] \le \EE[\EE[X^2 \mid \B]] = \EE[X^2] < \infty,
\]
and thus the random variable $\EE[X \mid \B]$ belongs to $L^2(\Omega,\B,\PP)$.

On the other hand, for every bounded $\B$-measurable random variable $Z$,
\[
\EE[Z(X - \EE[X \mid \B])] = \EE[ZX] - \EE[Z\EE[X \mid \B]] = 0,
\]
by the characteristic property (2) in Theorem \ref{thm:Conditional_expectation_existence} . Hence $X - \EE[X \mid \B]$ is orthogonal to the space of bounded $\B$-measurable random variables, and the latter space is dense in $L^2(\Omega,\B,\PP)$ (for instance by Lemma \ref{lem:approximation_simple_random_variables}). It follows that $X - \EE[X \mid \B]$ is orthogonal to $L^2(\Omega,\B,\PP)$, which gives the desired result.
}


\thmpr{Pull-out property of conditional expectation}{Pull_out_property}{
Let $X$ and $Y$ be real random variables and assume that $Y$ is $\B$-measurable. Then
\[ \EE[YX \mid \B] = Y\,\EE[X \mid \B], \]
provided $X$ and $Y$ are both nonnegative, or $X$ and $YX$ are both integrable.
}{
Suppose that $X \ge 0$ and $Y \ge 0$. Then, for every nonnegative $\B$-measurable random variable $Z$, we have
\[ \EE[Z(Y\EE[X \mid \B])] = \EE[(ZY)\EE[X \mid \B]] = \EE[ZYX], \]
using \textit{characteristic property} (3) from theorem \ref{thm:Conditional_expectation_existence} and the fact that $ZY$ is $\B$-measurable. Since $Y\EE[X \mid \B]$ is a nonnegative $\B$-measurable random variable, the characteristic property from theorem \ref{thm:Conditional_expectation_existence} implies that
\[ Y\EE[X \mid \B] = \EE[YX \mid \B]. \]

In the case when $X$ and $YX$ are integrable, we get the desired result by decomposing
\[ X = X^+ - X^-, \qquad Y = Y^+ - Y^- . \]
}

\thmpr{Nested $\sigma$-Fields}{conditional_nested_sigma_fields}{
Let $\B_1$ and $\B_2$ be two sub-$\sigma$-fields of $\A$ such that $\B_1 \subset \B_2$. Then, for every nonnegative (or integrable) random variable $X$, we have
\[ \EE[\EE[X \mid \B_2] \mid \B_1] = \EE[X \mid \B_1]. \]
}{
Consider the case where $X \ge 0$. Let $Z$ be a nonnegative $\B_1$-measurable random variable. Then, since $Z$ is also $\B_2$-measurable,
\[ \EE[Z\,\EE[\EE[X \mid \B_2] \mid \B_1]] = \EE[Z\,\EE[X \mid \B_2]] = \EE[ZX]. \]
Hence $\EE[\EE[X \mid \B_2] \mid \B_1]$ satisfies the characteristic property of $\EE[X \mid \B_1]$, and the desired result follows.
}
\textbf{Remark:} We have also $\EE[\EE[X \mid \B_1] \mid \B_2] = \EE[X \mid \B_1]$ under the same assumptions, but this is trivial since $\EE[X \mid \B_1]$ is $\B_2$-measurable.

\thmp{Independence and conditional expectaiton}{
Two sub-$\sigma$-fields $\B_1$ and $\B_2$ of $\A$ are independent if and only if, for every $B \in \B_2$, we have
\[ \EE[\1_B \mid \B_1] = \PP(B). \]
Furthermore, if $\B_1$ and $\B_2$ are independent, we have also $\EE[X \mid \B_1] = \EE[X]$ for every nonnegative $\B_2$-measurable random variable $X$ and for every $X \in L^1(\Omega,\B_2,\PP)$.
}{
Left exercise for the reader.
}

\thmpr{Factorization under independence}{factorization_under_independence_property}{
Let $(E,\E)$ and $(F,\F)$ be two measurable spaces, and let $X$ and $Y$ be random variables taking values in $E$ and $F$ respectively. Assume that $X$ is independent of $\B$ and that $Y$ is $\B$-measurable. Then, for every $\E \otimes \F$-measurable function $g : E \times F \to \RR_+$,
\[ \EE[g(X,Y) \mid \B] = \int g(x,Y)\,\PP_X(dx), \]
where we recall that $\PP_X$ denotes the law of $X$. The right-hand side is the composition of the random variable $Y$ with the mapping $\Psi : F \to \RR_+$ defined by
\[ \Psi(y) = \int g(x,y)\,\PP_X(dx). \]
}{
We have to prove that, for any nonnegative $\B$-measurable random variable $Z$, we have
\[ \EE[g(X,Y)Z] = \EE[\Psi(Y)Z]. \]

Write $\PP_{X,Y,Z}$ for the law of the triple $(X,Y,Z)$, which is a probability measure on $E \times F \times \RR_+$. Since $X$ is independent of $\B$, $X$ is independent of the pair $(Y,Z)$, and therefore
\[ \PP_{X,Y,Z} = \PP_X \otimes \PP_{Y,Z}. \]
Then, using the Fubini theorem, we have
\begin{align*}
\EE[g(X,Y)Z]
&= \int g(x,y)z \, \PP_{X,Y,Z}(dx\,dy\,dz) \\
&= \int g(x,y)z \, \PP_X(dx)\,\PP_{Y,Z}(dy\,dz) \\
&= \int_{F \times \RR_+} z \left( \int_E g(x,y)\,\PP_X(dx) \right) \PP_{Y,Z}(dy\,dz) \\
&= \int_{F \times \RR_+} z\,\Psi(y)\,\PP_{Y,Z}(dy\,dz) \\
&= \EE[\Psi(Y)Z].
\end{align*}
which was the desired result. 
}
\textbf{Remark:}
The mapping $\Psi$ is measurable thanks to the Fubini theorem. The theorem can be explained informally as follows. If we condition with respect to the $\sigma$-field $\B$, the random variable $Y$, which is $\B$-measurable, behaves like a constant, but on the other hand $\B$ gives no information on the random variable $X$. Thus the best approximation of $g(X,Y)$ knowing $\B$ is obtained by integrating $g(\cdot,Y)$ with respect to the law of $X$.

\thmpr{Invariance under independent information}{invariance_under_independent_information_property}{
Let $Z$ be a random variable in $L^1$, and let $\H_1$ and $\H_2$ be two sub-$\sigma$-fields of $\A$. Assume that $\H_2$ is independent of $\sigma(Z) \vee \H_1$. Then,
\[\EE[Z \mid \H_1 \vee \H_2] = \EE[Z \mid \H_1].\]
}{
It suffices to prove that the equality
\[ \EE[\1_A Z] = \EE[\1_A \EE[Z \mid \H_1]]  \]
holds.
}

\subsection{Transition Probabilities and Conditional Distributions}
\defn{Transition probability}{
Let $(E,\E)$ and $(F,\F)$ be two measurable spaces. A transition probability (or transition kernel) from $E$ into $F$ is a mapping
\[ \nu : E \times \F \longrightarrow [0,1] \]
which satisfies the following two properties:
\begin{enumerate}
\item for every $x \in E$, $A \mapsto \nu(x,A)$ is a probability measure on $(F,\F)$;
\item for every $A \in \F$, the function $x \mapsto \nu(x,A)$ is $\E$-measurable.
\end{enumerate}
}

At an intuitive level, each time one fixes a "starting point" $x \in E$, the probability measure $\nu(x,\cdot)$ gives a way to choose an "arrival point" $y \in F$ in a random manner (but with a law depending on the starting point $x$). This notion plays a fundamental role in the theory of Markov chains.

\propp{
Let $\nu$ be a transition probability from $E$ into $F$.
\begin{enumerate}
\item If $h$ is a nonnegative (resp.\ bounded) measurable function on $(F,\F)$, then
\[ \varphi(x) = \int h(y)\,\nu(x,dy), \qquad x \in E, \]
is a nonnegative (resp.\ bounded) measurable function on $E$.
\item If $\gamma$ is a probability measure on $(E,\E)$, then
\[\mu(A) = \int \gamma(dx)\,\nu(x,A), \qquad A \in \F,\]
is a probability measure on $(F,\F)$.
\end{enumerate}
}{
We omit the easy proof of the proposition. When treating the nonnegative case of (1), we first consider simple functions, and then use an increasing passage to the limit.  }

We now come to the relation between the notion of a transition probability and conditional expectations.

\defn{The conditional distribution}{
Let $X$ and $Y$ be two random variables taking values in $(E,\E)$ and $(F,\F)$ respectively. Any transition probability $\nu$ from $E$ into $F$ such that, for every nonnegative measurable function $h$ on $F$,
\[ \EE[h(Y) \mid X] = \int h(y)\,\nu(X,dy), \qquad \text{a.s.}, \]
is called a conditional distribution of $Y$ knowing $X$.
}

\textbf{Remark}
The random variable $\int h(y)\,\nu(X,dy)$ is the composition of $X$ with the function
\[ x \mapsto \int h(y)\,\nu(x,dy), \]
which is measurable by last Proposition. In particular, this random variable is a function of $X$, as $\EE[h(Y) \mid X]$ should be.

By definition, if $\nu$ is a conditional distribution of $Y$ knowing $X$, we have, for every $A \in \F$,
\[ \PP(Y \in A \mid X) = \nu(X,A), \qquad \text{a.s.} \]

It is tempting to replace this equality of random variables by an equality of real numbers,
\[ \PP(Y \in A \mid X = x) = \nu(x,A), \]
for every $x \in E$. Although it gives the intuition behind the notion of conditional distribution, this last equality has no mathematical meaning in general, because one has typically $\PP(X = x) = 0$ for every $x \in E$, which makes it impossible to define the conditional probability given the event $\{X = x\}$. The only rigorous formulation is thus the first equality $\PP(Y \in A \mid X) = \nu(X,A)$.

Let us briefly discuss the uniqueness of the conditional distribution of $Y$ knowing $X$. If $\nu$ and $\nu'$ are two such conditional distributions, we have, for every $A \in \F$,
\[ \nu(X,A) = \PP(Y \in A \mid X) = \nu'(X,A), \qquad \text{a.s.} \]
which is equivalent to saying that, for every $A \in \F$,
\[\nu(x,A) = \nu'(x,A), \qquad \PP_X(dx)\text{-a.s.}\]

Suppose that the measurable space $(F,\F)$ is such that any probability measure on $(F,\F)$ is characterized by its values on a (given) countable collection of measurable sets (this property holds for $(\RR^d,\B(\RR^d))$, by considering the collection of all rectangles with rational coordinates). Then, we conclude from the preceding display that
\[\nu(x,\cdot) = \nu'(x,\cdot), \qquad \PP_X(dx)\text{-a.s.}\]

So uniqueness holds in this sense (and clearly we cannot expect more). By abuse of language, we will nonetheless speak about \emph{the conditional distribution} of $Y$ knowing $X$.

Let us now consider the problem of the existence of conditional distributions.

\thmp{Existence of conditional distribution}{
Let $X$ and $Y$ be two random variables taking values in $(E,\E)$ and $(F,\F)$, respectively. Suppose that $(F,\F)$ is a complete separable metric space equipped with its Borel $\sigma$-field. Then, there exists a conditional distribution of $Y$ knowing $X$.
}{
We will not prove this theorem, see Theorem 6.3 in \cite{Probability2} for the proof of a more general statement.
}

\subsection{Evaluation of Conditional Expectation}
\subsubsection{Discret Case}
If $X$ is a discrete variable ($E$ is countable), and $Y$ any random variable. then we have the conditional distibution of $Y$ given $X$, denoted by $\nu(x,A)$, defined by :
\[ \nu(x,A) := 
\begin{cases} 
    \frac{1}{\PP(X=x)}\PP(Y \in A, X = x) & \text{if } x \in E' := \{y \in E : \PP(X = y) > 0\}, \\[6pt]
    \delta_{y_0}(A) & \text{if } x \notin E', 
\end{cases} \]
where $y_0$ is an arbitrary fixed point of $F$. Theorefore we can conclude that : 
\[ \EE[h(Y) \mid X] = \varphi(X) \]
where
\[ \varphi(x) = \int h(y)\,\nu(x,dy) = 
\begin{cases} 
    \frac{1}{\PP(X=x)}\EE(h(Y)\1_{\{X=x\}}) & \text{if } x \in E' , \\[6pt]
    y_0 & \text{if } x \notin E', 
\end{cases} \]
Where $y_0$ is an arbitrary fixed point of $F$.

\subsubsection{Continuous Case [Random Variables with a Density]}
Suppose that $X$ and $Y$ take values in $\RR^m$ and in $\RR^n$ respectively, and that the pair $(X,Y)$ has density $p(x,y)$, $(x,y) \in \RR^m \times \RR^n$. The density of $X$ is then
\[q(x) = \int_{\RR^n} p(x,y)\, dy,\]
where we take $q(x) = 0$ if the integral is infinite. We may define the conditional distribution of $Y$ knowing $X$ by :
\[ \nu(x,A) = 
\begin{cases} 
    \frac{1}{q(x)} \int_A p(x,y)\, dy & \text{if } q(x) > 0, \\[6pt]
    \delta_0(A) & \text{if } q(x) = 0. 
\end{cases} 
\]
 Theorefore we can conclude that : 
\[ \EE[h(Y) \mid X] = \varphi(X) \]
where
\[ \varphi(x) = \int h(y)\,\nu(x,dy) = 
\begin{cases}
    \dfrac{1}{q(y)} \displaystyle \int_{\RR^m} h(x)\, p(x,y)\, dx & \text{if } q(y) > 0, \\
    h(0) & \text{if } q(y) = 0.
\end{cases}\]



\section{Convergence of random variables} 

The study of convergence of random variables is fundamental in probability theory, particularly in the analysis of asymptotic behaviors and limit theorems. Different notions of convergence allow us to formalize how a sequence of random variables approximates another random variable under various probabilistic criteria. These modes of convergence differ in strength and applicability, and understanding their relationships is essential for results such as the law of large numbers and the central limit theorem.

Throughout this section, we consider a probability space $(\probspace)$. Let $(\bX_n)_{n\in\NN}$ be a sequence of random variables that mapsto $\RR^d$ and let $\bX$ be another random variable taking values in $\RR^d$. We introduce below the principal notions of convergence that will be used in the sequel.

\defnr{Almost sure convergence}{as_conv}{
The sequence $(\bX_n)$ is said to converge almost surely to $\bX$, denoted by
\[\bX_n \conv[a.s.]{n} \bX, \]
if
\[\PP\left(\left\{\omega \in \Omega :\lim_{n\to\infty} \bX_n(\omega)=\bX(\omega)\right\}\right)=1.\]}

The \emph{Borel--Cantelli Lemma} \ref{thm:Borel_Cantelli} is a powerful tool for proving (or disproving) almost sure convergence. The main idea is to convert the convergence problem into the study of a sequence of \textit{bad}  events as discribed in Figure \ref{fig:borel_cantelli_strategy} and define the bad events as follow :
\[A_n=\{|X_n-X|>\varepsilon\}, \]


\defnr{Convergence in probability}{prob_cong}{
We say that the sequence $(\bX_n)_{n\in\NN}$ converges in probability to $\bX$, and write
\[\bX_n \conv[\PP]{n} \bX,\]
if for every $\varepsilon>0$,
\[\lim_{n\to\infty}\PP\big(|\bX_n - \bX|>\varepsilon\big)=0.\]}


\proppr{prob_conv_characterisation}{
We remind that $d$ is the distance on $L^0(\Omega,\A,\PP)$ defined by :
\[ d(\bX,\bY) = \EE[|\bX-\bY| \wedge  1]\] 
We have : 
\[ \bX_n \conv[\PP]{n} \bX \Longleftrightarrow d(\bX_n,\bX) \conv{n} 0\]
}{
Left exercice for the reader. \\  
Hint: use the fact that the space $L^0_{\RR^d}(\Omega,\A,\PP)$ is complete for the distance $d$ by proposition \ref{prop:L0_complete}.
}

\defnr{Convergence in Lp}{Lp_conv} {
Let $p\in[1,\infty)$ and assume that
\[\EE[|\bX|^p]<\infty\quad \text{and} \quad\EE[|\bX_n|^p]<\infty \quad \text{for all } n\in\NN.\]
The sequence $(\bX_n)$ is said to converge to $\bX$ in $L^p$, written
\[\bX_n \conv[\L^p]{n} \bX,\]
if
\[\lim_{n\to\infty}\EE\big[|\bX_n - \bX|^p\big]=0.\]}
This is equivalent to saying that each component sequence of $(\bX_n)_{n \in \NN}$ converges to the corresponding component of $\bX$ in the Banach space $L^p(\Omega,\A,\PP)$. One could also consider the convergence in $L^\infty$, but this convergence is very rarely used in probability theory and we will not discuss it.

A relevant remark is that domination transforms a pointwise convergence problem into an \(L^p\)-convergence considerably more straightforward.

\subsection{Relationship between modes of convergences }
The different notions of convergence introduced in this chapter are related by a number of fundamental implications. Among them, convergence in probability occupies an intermediate position: it is weaker than both almost sure convergence and $L^p$-convergence, yet stronger than convergence in distribution. Consequently, every almost sure convergence or $L^p$-convergence yields convergence in probability, while convergence in probability itself implies convergence in distribution.

The latter notion is introduced and studied in the following sections. Briefly, convergence in distribution describes the convergence of the laws $(\PP_{X_n})_{n\in\NN}$ toward the law $\PP_X$.

The three direct implications are summarized by the following chain:
\begin{itemize}
    \item almost sure convergence imply probability
    \item $L^p$-convergence imply probability
    \item convergence in probability imply convergence in distribution
\end{itemize}
The converse implications are, in general, false. Throughout this section, we establish the direct implications above and investigate the conditions under which some of the converse implications may nevertheless hold.
\proppr{prob_imply_dist}{
    If  
    \[X_n\conv[\PP]{n} X\] 
    Then : 
    \[X_n\conv[\D]{n} X\] 
}{
Suppose first that $X_n$ converges a.s. to $X$. Then, for every $\varphi\in C_b(\RR^d)$, $\varphi(X_n)$ converges a.s. to $\varphi(X)$ and the dominated convergence theorem gives $\EE[\varphi(X_n)] \to \EE[\varphi(X)]$.

If $X_n$ only converges in probability to $X$, then we know from Proposition \ref{prop:prob_imply_subseq_as} that there is a subsequence of $(X_n)_{n\in\NN}$ that converges a.s. to $X$, and thus, for every $\varphi\in C_b(\RR^d)$, $\EE[\varphi(X_n)]$ converges to $\EE[\varphi(X)]$ along this subsequence. But then, the same argument shows that, from any subsequence of $(X_n)_{n\in\NN}$, we can extract a subsequence along which $\EE[\varphi(X_n)]$ converges to $\EE[\varphi(X)]$. This is only possible if $\EE[\varphi(X_n)] \to \EE[\varphi(X)]$ as $n\to\infty$. }

\textbf{Remark:} The converse to the proposition is (of course) false, because, as we already mentioned, the convergence in distribution of $X_n$ to $X$ does not determine the limiting variable $X$. 

\proppr{as_or_Lp_imply_Prob}{
    If : 
    \[ 
        \bX_n \conv[a.s.]{n} \bX, \quad \text{or} \quad \bX_n \conv[\L^p]{n} \bX.
    \]
    Then : 
    \[ 
        \bX_n \conv[\PP]{n} \bX.
    \]
}{
\begin{enumerate}
    \item If $\bX_n$ converges to $\bX$ in almost surely, then define :
    $$ Z_n = \1_{\{||\bX_n-\bX||>\varepsilon\}} $$
    Since $||\bX_n-\bX|| \conv{n}  0$ then $Z_n$ converge almost surely to $0$, then apply the Dominated Convergence Theorem to conclude that :
    $$\PP(||\bX_n-\bX||>\varepsilon) = \EE[ \1_{\{||\bX_n-\bX||>\varepsilon\}}] = \EE[Z_n] \conv{n}  0$$
    And therefore the convergence in probability. 
    \item If $\bX_n$ converges to $\bX$ in $L^p$, then by Markov inequality \ref{thm:Markov_Inequalities} we conclude the convergence in probability as follow : 
    $$P(||X_n-X||>\varepsilon) =  P(||X_n-X||^p>\varepsilon^p) \le \frac{E|X_n-X|^p}{\varepsilon^p} \conv{n} 0$$ 
\end{enumerate}
}
Although convergence in probability is weaker than almost sure convergence, it has a key compactness property: every sequence converging in probability admits a subsequence that converges almost surely to the same limit. 
\proppr{prob_imply_subseq_as}{
    If : 
    \[ \bX_n \conv[\PP]{n} \bX. \]
    Then : 
    \[ \exists, (n_k) \uparrow \infty \quad \text{ such that }\quad  \bX_{n_k} \conv[a.s.]{k} \bX \]
}{
Choose the subsequence so that
$$P(\vert{}X_{n_k} - X\vert{} > 2^{-k}) < 2^{-k}.$$
Then
$$\sum_{k} P(\vert{}X_{n_k} - X\vert{} > 2^{-k}) < \infty.$$
Apply the first Borel-Cantelli lemma.
Eventually,
$$\vert{}X_{n_k} - X\vert{} \le 2^{-k},$$
hence almost sure convergence.
}

The previous result only yields almost sure convergence along a subsequence and cannot, in general, be strengthened to full almost sure convergence. One may then ask whether convergence in probability can be upgraded to ($L^p$)-convergence. This is false without extra assumptions. A key sufficient condition is a uniform ($L^r$) bound with ($r>p$), which ensures ($L^p$)-convergence via uniform integrability.
\proppr{prob_imply_Lp}{
    If : 
    \[ 
    \begin{cases*}
        \bX_n \conv[\PP]{n} \bX. \\
        \exists r \in (1,\infty) \quad \text{ such : }\quad  (\EE[|\bX_n|^r])_{n \in \NN} \quad \text{is bounded}.
    \end{cases*}
    \]
    Then : 
    \[ 
    \begin{cases*}
        \EE[|\bX|^r] < \infty \\ 
        \forall p \in [1,r), \quad \bX_n \conv[\L^p]{n} \bX.
    \end{cases*}
    \]
}{
By assumption, there is a constant $C$ such that $\EE[|\bX_n|^r] \le C$ for every $n$. By the previous proposition we can find a subsequence $(\bX_{n_k})$ that converges almost surely to $\bX$. By Fatou's lemma,

\begin{align*}
\EE[|\bX|^r]
&= \EE\Bigl[\liminf_{k \to \infty} |\bX_{n_k}|^r\Bigr] \\
&\le \liminf_{k \to \infty} \EE[|\bX_{n_k}|^r] \\
&\le C.
\end{align*}

Then, for every $p \in [1,r)$ and every $\varepsilon > 0$,
\begin{align*}
\EE[|\bX_n - \bX|^p] 
    &= \EE\bigl[|\bX_n - \bX|^p \1_{\{|\bX_n - \bX| \le \varepsilon\}}\bigr]  + \EE\bigl[|\bX_n - \bX|^p \1_{\{|\bX_n - \bX| > \varepsilon\}}\bigr] \\
    &\le \varepsilon^p + \EE[|\bX_n - \bX|^r]^{p/r}\PP(|\bX_n - \bX| > \varepsilon)^{1 - p/r} \\
    &\le \varepsilon^p + (2C)^{p/r} \PP(|\bX_n - \bX| > \varepsilon)^{1 - p/r}.
\end{align*}

Since $\bX_n$ converges in probability to $\bX$,
\[\limsup_{n \to \infty} \EE[|\bX_n - \bX|^p] \le \varepsilon^p.\]
Letting $\varepsilon \to 0$ yields $\EE[|\bX_n - \bX|^p] \to 0$.
}

Another important criterion for strengthening convergence in probability is provided by Scheffé's lemma. Instead of requiring higher-order moment bounds, it assumes the convergence of the first moments. Under the additional assumptions of non-negativity and convergence of expectations, convergence in probability is sufficient to conclude convergence in ($L^1$).
\lempr{Scheffé's lemma}{prob_imply_L1}{
    If : 
    \[ 
    \begin{cases*}
        \bX_n \geq 0 \quad \text{and}\quad  \bX_n \conv[\PP]{n} \bX. \\
        \forall n\in \NN :\quad \bX_n\in L^1 \quad \text{and}\quad \bX \in L^1 \\
        \EE[\bX_n] \conv{n} \EE[\bX]\\
    \end{cases*}
    \]
    Then : 
    \[ 
        \bX_n \conv[\L^1]{n} \bX.
    \]
}{
We write
\[ \EE[|\bX_n - \bX|] = \EE[\bX_n - \bX] + 2\EE[(\bX_n - \bX)^-]. \]
We have $\EE[\bX_n - \bX] \to 0$, and the bound $(\bX_n - \bX)^- \le \bX$ allows us to apply the dominated convergence theorem to obtain $\EE[(\bX_n - \bX)^-] \to 0$ (using that almost sure convergence in the dominated convergence theorem can be replaced by convergence in probability).
}
A more general criterion for upgrading convergence in probability to ($L^1$)-convergence is provided by \textbf{Vitali's theorem}, which replaces moment assumptions by the weaker notion of uniform integrability. Since the proof of this result relies on techniques from martingale theory, it will be deferred to the next chapter. We therefore conclude this section by introducing the notion of uniform integrability.

A family $(X_i)_{i\in I}$ of random variables in $L^1(\Omega,\A,\PP)$ is said to be \textit{uniformly integrable (u.i.)} if
\[ \lim*{a\to+\infty} \sup_{i\in I} \EE!\left[|X_i|\1_{{|X_i|>a}}\right] =0. \]

\thmpr{Vitali convergence theorem}{prob_imply_L1_vitali}{
If 
\[ \begin{cases*}
    \bX_n \conv[\PP]{n} \bX. \\
    \bX_n \in L^1 \quad \text{for every } n\in\NN.\\
    \text{the sequence } (\bX_n)_{n \in \NN} \text{ is u.i.} 
\end{cases*} \] 
Then the following assertions are equivalent :
\[
    \bX_n \conv[L^1]{n} \bX  
\]
}{
    It would be shown in the section of uniform integrability \ref{sec:uniform_integrability} (next chapter).  
}
It's also relevant to see that if there is $Z\in L^1$ such $|\bX_n| \leq Z$ a.s. for every $n\in\NN$, then the sequence $(\bX_n)_{n \in \NN}$ is u.i., which gives a simpler condition to prove $L^1$-convergence. In fact, we have :
\[ \EE[|\bX_n \1_{|\bX_n\geq a}] \leq \EE[Z\1_{Z\geq a}]  \]
Finaly, taking supremum and limits : 
\[ \lim_{\goto{a}} \sup_{n\in\NN} \EE[|\bX_n| \1_{|\bX_n\geq a}] \leq \lim_{\goto{a}} \EE[Z\1_{Z\geq a}] = 0 \quad \text{Because :}\quad Z\1_{Z\geq a} \conv[a.s.]{n} 0\]



\begin{figure}[h]
\centering
\[
\begin{tikzcd}[column sep=huge, row sep=large]
|[draw, rounded corners, inner sep=3pt]| {X_n \xrightarrow{L^q} X}
  \arrow[r, "q \ge p \ge 1"', Rightarrow]
&
|[draw, rounded corners, inner sep=3pt]| {X_n \xrightarrow{L^p} X}
  \arrow[r, Rightarrow]
&
|[draw, rounded corners, inner sep=3pt]| {X_n \xrightarrow{L^1} X}
  \arrow[d, Rightarrow]
\\
&&
|[draw, rounded corners, inner sep=3pt]| {X_n \xrightarrow{\PP} X}
  \arrow[r, Rightarrow]
  \arrow[lu, "{L^r-Bounded \quad r>p}", dashed, bend left=35]
  \arrow[u, "{\text{Expectation convergence}}"', dashed, bend right=45]
  \arrow[u, swap, "{\substack{\text{Uniform conv.}\\ \text{or domination}}}"', dashed, bend left=45]
  \arrow[d, "{\text{along a subsequence}}", dashed, bend left=60]
&
|[draw, rounded corners, inner sep=3pt]| {X_n \xrightarrow{\mathcal D} X}
\\
&&
|[draw, rounded corners, inner sep=3pt]| {X_n \xrightarrow{\text{a.s.}} X}
  \arrow[u, Rightarrow]
\end{tikzcd}
\]
\caption{Relations between the different convergences} \label{fig:relationshipsConv}
\end{figure}


Figure \ref{fig:relationshipsConv} illustrates the relations between the different types of convergence that we have introduced. The dashed lines give partial converses (under the stated conditions) that correspond to results stated in Propositions \ref{prop:as_or_Lp_imply_Prob}, \ref{prop:prob_imply_subseq_as}, \ref{lem:prob_imply_L1}, \ref{prop:prob_imply_Lp}, and to the dominated convergence theorem.


\subsection{Convergence in Distribution}
Recall that $C_b(\RR^d)$ denote the set of all bounded continuous functions from $\RR^d$ into $\RR$. We equip $C_b(\RR^d)$ with the supremum norm 
\[ \|\varphi\|=\sup_{x\in\RR^d}|\varphi(x)|. \]
Also recall the notation $C_c(\RR^d)$ for the space of all real continuous functions with compact support on $\RR^d$.
We let $M_1(\RR^d)$ denote the set of all probability measures on $(\RR^d,\B(\RR^d))$.

\defnr{Convergence in distribution}{dist_conv}{A sequence $(\QQ_n)_{n\in\NN}$ in $M_1(\RR^d)$ converges weakly to $\QQ\in M_1(\RR^d)$ if
\[ \forall \varphi\in C_b(\RR^d), \qquad \int \varphi\, d\QQ_n \conv{n} \int \varphi\, d\QQ . \]
A sequence $(X_n)_{n\in\NN}$ of random variables with values in $\RR^d$ converges in distribution to a random variable $X$ with values in $\RR^d$ if the sequence $(P_{X_n})_{n\in\NN}$ converges weakly to $P_X$.\\
By Properties of Expectation \ref{defn:Expectation}, this is equivalent to saying that
\[\forall \varphi\in C_b(\RR^d), \qquad \EE[\varphi(X_n)] \conv{n} \EE[\varphi(X)].\]

We will write $\QQ_n \conv[(w)]{n} \QQ$, resp. $X_n \conv[\D]{n} X$, if the sequence $(\QQ_n)_{n\in\NN}$ converges weakly to $\QQ$, resp. if $(X_n)_{n\in\NN}$ converges in distribution to $X$.
}

\textbf{Remark}
\begin{itemize}
\item There is some abuse of terminology in saying that the sequence $(X_n)_{n\in\NN}$ converges in distribution to $X$, because the limiting random variable $X$ is not determined uniquely (only its distribution $P_X$ is determined). For this reason, we will sometimes write that a sequence $(X_n)_{n\in\NN}$ of random variables converges in distribution to a probability measure $\QQ$ --- one should of course understand that the laws $P_{X_n}$ converge weakly to $\QQ$. We may also note that the convergence in distribution makes sense even if the random variables $X_n$, $n\in\NN$ are defined on different probability spaces (in the present book, we will however always assume that they are defined on the same probability space). This makes the convergence in distribution very different from the other types of convergence studied in this section.
\item Although the definition of weak convergence is formulated in terms of all bounded continuous functions, such a family is unnecessarily large in practice. It is therefore natural to ask whether convergence only needs to be verified on a smaller collection of test functions. The next proposition shows that this is indeed the case.
\end{itemize}

\proppr{weaker_families_on_weak_conv}{Let $(\QQ_n)_{n\in\NN}$ and $\QQ$ be probability measures on $\RR^d$. Let $H$ be a subset of $C_b(\RR^d)$ whose closure (with respect to the supremum norm) contains $C_c(\RR^d)$. The following properties are equivalent.
\begin{enumerate}
\item  The sequence $(\QQ_n)_{n\in\NN}$ converges weakly to $\QQ$.
\item We have
\[\forall \varphi \in C_c(\RR^d), \qquad \int \varphi\, d\QQ_n \conv{n} \int \varphi\, d\QQ .\]
\item  We have
\[\forall \varphi \in H, \qquad \int \varphi\, d\QQ_n \conv{n} \int \varphi\, d\QQ .\]
\end{enumerate}}{
It is obvious that (1)$\Rightarrow$(2) and (1)$\Rightarrow$(3). Then suppose that (2) holds. Let
$\varphi \in C_b(\RR^d)$ and let $(f_k)_{k\in\NN}$ be a sequence in $C_c(\RR^d)$ such that
$0 \le f_k \le 1$ for every $k \in \NN$, and $f_k \uparrow 1$ as $k \to \infty$. Then, for every $k \in \NN$,
$\varphi f_k \in C_c(\RR^d)$ and thus
\[ \int \varphi f_k \, d\QQ_n \conv{n} \int \varphi f_k \, d\QQ. \tag{i} \]

On the other hand,
\[ \left| \int \varphi \, d\QQ_n - \int \varphi f_k \, d\QQ_n \right| \le \|\varphi\| \int (1-f_k)\, d\QQ_n = \|\varphi\|\left(1-\int f_k\, d\QQ_n\right), \]

and similarly

\begin{align*}
\left| \int \varphi \, d\QQ - \int \varphi f_k \, d\QQ \right|
&\le \|\varphi\|\left(1-\int f_k\, d\QQ\right).
\end{align*}

Hence, using (i), we get for every $k \in \NN$,
\begin{align*}
\limsup_{n\to\infty}
\left| \int \varphi\, d\QQ_n - \int \varphi\, d\QQ \right|
    &\le \limsup_{n\to\infty}
            \left|  \int \varphi\, d\QQ_n - \int \varphi f_k\, d\QQ_n\right|
            + \left| \int \varphi\, d\QQ - \int \varphi f_k\, d\QQ\right| \\
    &\le \|\varphi\|
        \left( \limsup_{n\to\infty} \left(1-\int f_k\, d\QQ_n\right) + \left(1-\int f_k\, d\QQ\right) \right) \\
    &=2\|\varphi\| \left(1-\int f_k\, d\QQ\right).
    \end{align*}

Now we just have to let $k$ tend to $\infty$ (noting that $\int f_k\, d\QQ \uparrow 1$, by monotone convergence) and we find that $\int \varphi\, d\QQ_n$ converges to $\int \varphi\, d\QQ$, so that (1) holds.

We complete the proof by verifying that (3)$\Rightarrow$(2). So assume that (3) holds. If $\varphi \in C_c(\RR^d)$, then, for every $k \in \NN$, we can find a function $\varphi_k \in H$ such that
$\|\varphi-\varphi_k\|\le 1/k$, and it follows that

\begin{align*}
    & \limsup _{n \rightarrow \infty}\left|\int \varphi \mathrm{~d} \QQ_n-\int \varphi \mathrm{d} \QQ\right| \\
    & \leq \limsup _{n \rightarrow \infty}\left(\left|\int\left(\varphi-\varphi_k\right) \mathrm{d} \QQ_n\right|+\left|\int \varphi_k \mathrm{~d} \QQ_n-\int \varphi_k \mathrm{~d} \QQ\right|+\left|\int\left(\varphi_k-\varphi\right) \mathrm{d} \QQ\right|\right) \\
    & \leq \frac{2}{k}
\end{align*}
Since this holds for every $k$, we get that
\[\int \varphi\, d\QQ_n \longrightarrow \int \varphi\, d\QQ.\]
}

The previous proposition simplifies the family of test functions, but weak convergence is still expressed through integrals against continuous functions. It is natural to ask whether one can characterize weak convergence without referring to test functions at all. The celebrated Portmanteau Theorem provides several equivalent descriptions in purely measure-theoretic terms, involving only the behavior of probability measures on open, closed, or continuity sets.

\thmpr{Portmanteau Theorem}{Portmanteau_theorem}{Let $(\QQ_n)_{n\in\NN}$ be a sequence in $M_1(\RR^d)$ and let $\QQ\in M_1(\RR^d)$. The following four assertions are equivalent:
\begin{enumerate}
    \item  The sequence $(\QQ_n)$ converges weakly to $\QQ$.
    \item  For every open subset $G$ of $\RR^d$,
    \[\liminf \QQ_n(G) \ge \QQ(G).\]
    \item  For every closed subset $F$ of $\RR^d$,
    \[\limsup \QQ_n(F) \le \QQ(F).\]
    \item  For every Borel subset $B$ of $\RR^d$ such that $\QQ(\partial B)=0$,
    \[\lim \QQ_n(B) = \QQ(B).\]
\end{enumerate}

}{
Let us start by proving that (1)$\Rightarrow$(2). If $G$ is an open subset of $\RR^d$, we can find a sequence $(\varphi_p)_{p\in\NN}$ of bounded continuous functions such that $0 \le \varphi_p \le \1_G$ and $\varphi_p \uparrow \1_G$ (for instance, $\varphi_p(x)=p\,\mathrm{dist}(x,G^c)\wedge 1$). Then
\[ \liminf_{n\to\infty} \QQ_n(G) \ge \sup_p \left(\liminf_{n\to\infty} \int \varphi_p\, d\QQ_n\right) = \sup_p \left(\int \varphi_p\, d\QQ\right) = \QQ(G). \]

The equivalence (2)$\Leftrightarrow$(3) is immediate since complements of open sets are closed and conversely.

Let us show that (2) and (3) imply (4). If $B\in\B(\RR^d)$, write $\overline{B}$ for the closure of $B$ and $B^\circ$ for the interior of $B$ so that $\partial B=\overline{B}\setminus B^\circ$. By (2) and (3), we have : 
\[\begin{cases}
    \limsup \QQ_n(B) \le& \limsup \QQ_n(\overline{B}) \le \QQ(\overline{B})\\
    \liminf \QQ_n(B) \ge& \liminf \QQ_n(B^\circ) \ge \QQ(B^\circ).
\end{cases}\]

If $\QQ(\partial B)=0$ then $\QQ(\overline{B})=\QQ(B^\circ)=\QQ(B)$ and we get (4).

We still have to prove that (4)$\Rightarrow$(1). Let $\varphi\in C_b(\RR^d)$. Without loss of generality, we can assume that $\varphi \ge 0$. Then, let $K \ge 0$ such that $0 \le \varphi \le K$. By the Fubini theorem,
\[\int \varphi(x)\QQ(dx)= \int \left(\int_0^K \1_{\{t\le\varphi(x)\}} dt \right)\QQ(dx)= \int_0^K \QQ(E_t^\varphi)\,dt,\]
where $E_t^\varphi$ is the closed set $E_t^\varphi := \{x\in\RR^d : \varphi(x)\ge t\}$. Similarly, for every $n$,
\[ \int \varphi(x)\QQ_n(dx) = \int_0^K \QQ_n(E_t^\varphi)\,dt. \]

Now note that $\partial E_t^\varphi \subset \{x\in\RR^d : \varphi(x)=t\}$, and that there are at most countably many values of $t$ such that
\[ \QQ(\{x\in\RR^d : \varphi(x)=t\}) > 0 \]

(indeed, for every $k\in\NN$, there are at most $k$ distinct values of $t$ such that $\QQ(\{x\in\RR^d : \varphi(x)=t\}) \ge \frac{1}{k}$, since the sets $\{x\in\RR^d : \varphi(x)=t\}$ are disjoint when $t$ varies, and $\QQ$ is a probability measure). Hence (4) implies
\[ \QQ_n(E_t^\varphi) \conv{n} \QQ(E_t^\varphi), \]

for every $t\in[0,K]$, except possibly for a countable set of values of $t$, which has zero Lebesgue measure. By dominated convergence, we conclude that
\[ \int \varphi(x)\QQ_n(dx) = \int_0^K \QQ_n(E_t^\varphi)\,dt \conv{n} \int_0^K \QQ(E_t^\varphi)\,dt = \int \varphi(x)\QQ(dx). \]
}



Although the Portmanteau Theorem gives several equivalent formulations of weak convergence, these criteria are often difficult to verify in concrete problems. A much more tractable approach is obtained by passing to characteristic functions. Lévy's Continuity Theorem shows that weak convergence is completely characterized by pointwise convergence of Fourier transforms. Recall our notation $\hat{\QQ}(\xi) = \int e^{i\xi\cdot x}\,\QQ(dx)$ for the Fourier transform of $\QQ \in M_1(\RR^d)$, and $\Phi_X=\widehat{P_X}$ for the characteristic function of a random variable $X$
with values in $\RR^d$.

\thmpr{Lévy's Theorem}{Levy_theorem}{
A sequence $(\QQ_n)_{n\in\NN}$ in $M_1(\RR^d)$ converges weakly to $\QQ \in M_1(\RR^d)$ if and only if
\[\forall \xi \in \RR^d, \qquad \hat{\QQ}_n(\xi) \xrightarrow[n\to\infty]{} \hat{\QQ}(\xi).
\]

Equivalently, a sequence $(X_n)_{n\in\NN}$ of random variables with values in $\RR^d$ converges in distribution to $X$ if and only if
\[\forall \xi \in \RR^d, \qquad \Phi_{X_n}(\xi) \xrightarrow[n\to\infty]{} \Phi_X(\xi).
\]
}{
It is enough to prove the first assertion. First, if we assume that the sequence
$(\QQ_n)_{n\in\NN}$ converges weakly to $\QQ$, the definition of this convergence shows that
\[ \forall \xi \in \RR^d , \qquad \hat{\QQ}_n(\xi) = \int e^{i\xi\cdot x}\,\QQ_n(dx) \xrightarrow[n\to\infty]{} \int e^{i\xi\cdot x}\,\QQ(dx) = \hat{\QQ}(\xi). \]

Conversely, assume that $\hat{\QQ}_n(\xi)\to \hat{\QQ}(\xi)$ for every $\xi \in \RR^d$ and let us show that
$(\QQ_n)$ converges weakly to $\QQ$. To simplify notation, we only treat the case $d=1$, but
the proof in the general case is exactly similar.

Let $f \in C_c(\RR)$, and, for every $\sigma>0$, let
\[ g_\sigma(x)=\frac{1}{\sigma\sqrt{2\pi}}\exp\!\left(-\frac{x^2}{2\sigma^2}\right). \]

be the density of the $\N(0,\sigma^2)$ distribution. We already observed in the proof of
Theorem \ref{thm:law_determined_by_characteristic_func} that $g_\sigma * f$ converges pointwise to $f$ as $\sigma \to 0$. Since $f$ has compact
support here, it is not hard to verify that this convergence is uniform on $\RR$ (use
Proposition \ref{prop:dirac_limit} (i) and note that, for every $\varepsilon>0$, we can find a compact set $K$ such
that $|g_\sigma * f|\le \varepsilon$ for every $x\in\RR\setminus K$ and $\sigma\in(0,1]$).
Thus, if we let $H$ be the subset of $C_b(\RR^d)$ defined by
\[ H:=\{\varphi=g_\sigma * f:\sigma>0 \text{ and } f\in C_c(\RR^d)\} \]

the closure of $H$ in $C_b(\RR^d)$ contains $C_c(\RR^d)$.

On the other hand, we also saw in the proof of Theorem \ref{thm:law_determined_by_characteristic_func} that, for any $\nu\in M_1(\RR)$,
\begin{align*}
\int g_\sigma * f \mathrm{~d} v & =\int f(x) g_\sigma * v(x) \mathrm{d} x \\
& =\int f(x)\left((\sigma \sqrt{2 \pi})^{-1} \int e^{i \xi x} g_{1 / \sigma}(\xi) \widehat{v}(-\xi) \mathrm{d} \xi\right) \mathrm{d} x \tag{1}
\end{align*}

We apply this formula to $\nu=\QQ_n$ and $\nu=\QQ$. From our assumption
$\hat{\QQ}_n(\xi)\to\hat{\QQ}(\xi)$ for every $\xi\in\RR$ and to the trivial bound
$|\hat{\QQ}_n(\xi)|\le1$, we can apply the dominated convergence theorem to get that, for every
$x\in\RR$,
\[ \int e^{i\xi x} g_{1/\sigma}(\xi)\hat{\QQ}_n(-\xi)\,d\xi \xrightarrow[n\to\infty]{} \int e^{i\xi x} g_{1/\sigma}(\xi)\hat{\QQ}(-\xi)\,d\xi . \]

Since the quantities in the last display are bounded by $1$ and $f\in C_c(\RR)$, we can use (1) and again the dominated convergence theorem to get
\[ \int g_\sigma * f\, d\QQ_n \xrightarrow[n\to\infty]{} \int g_\sigma * f\, d\QQ . \]

Finally, we have proved that $\int \varphi\, d\QQ_n \to \int \varphi\, d\QQ$ for every
$\varphi\in H$, and we know that the closure of $H$ in $C_b(\RR^d)$ contains
$C_c(\RR^d)$. By Proposition \ref{prop:weaker_families_on_weak_conv}, this is enough to show that the sequence
$(\QQ_n)_{n\in\NN}$ converges weakly to $\QQ$. 
}




\subsection{The Strong Law of Large Numbers}
Let us begin with a useful implication of Kolmogorov's Zero--One Law (Theorem~\ref{thm:zero_one_law}). Suppose that $(X_n)_{n\ge 1}$ is a sequence of independent random variables and that
\[ Z=\lim_{n\to\infty} \frac{1}{n}(X_1+\ldots+X_n) \]
exists almost surely. For every fixed $k\in\NN$, the random variable $Z$ may equivalently be expressed as a measurable function depending only on the tail variables
\[ Z =  \lim_{n\to\infty} \frac{1}{n}(X_k+X_{k+1}+\ldots) \in \sigma(X_k, X_{k+1},\ldots)  \]
then
\[ Z\in\bigcap_{k=1}^{\infty}\sigma(X_k,X_{k+1},\ldots), \]
that is, $Z$ is measurable with respect to the tail $\sigma$-algebra. Consequently, by Theorem~\ref{thm:zero_one_law}, the random variable $Z$ is almost surely constant.


\thmpr{Strong Law of Large Numbers}{strong_law_large_numbers}{
Let $(X_n)_{n \ge 1}$ be a sequence of independent and identically distributed real random variables in $L^1$. Then
\[ \frac{1}{n}(X_1 + \cdots + X_n) \conv[a.s.]{n} \EE[X_1]. \]
}{
To simplify notation, we set $S_n = X_1 + \cdots + X_n$ for every $n \in \NN$, and $S_0 = 0$. Let $a > \EE[X_1]$, and
\[ M := \sup_{n \in \ZZ_+} (S_n - na), \]
which is a random variable with values in $[0,\infty]$. \\ 
We will show that
\[ M < \infty, \quad \text{a.s.} \tag{1} \]

Assume that (1) holds (for any choice of $a > \EE[X_1]$). \\ 
Since $S_n \le na + M$ for every $n \in \NN$, it immediately follows from (1) that
\[ \limsup_{n \to \infty} \frac{1}{n} S_n \le a, \quad \text{a.s.} \]
By considering a sequence of values of $a$ that decreases to $\EE[X_1]$, we find
\[ \limsup_{n \to \infty} \frac{1}{n} S_n \le \EE[X_1], \quad \text{a.s.} \]
Replacing $X_n$ by $-X_n$, we get the reverse inequality
\[ \liminf_{n \to \infty} \frac{1}{n} S_n \ge \EE[X_1], \quad \text{a.s.} \]
and the convergence in the theorem follows.

It remains to prove (1). With the notation of Theorem \ref{thm:zero_one_law}, we first observe that $\{M < \infty\}$ belongs to $\B_\infty$. Indeed, we have, for every integer $k \in \NN$,
\[ \{M < \infty\} = \left\{ \sup_{n \in \ZZ_+} (S_n - na) < \infty \right\} = \left\{ \sup_{n \ge k} (S_n - S_k - na) < \infty \right\}, \]
and the event in the right-hand side is $\B_{k+1}$-measurable. It will therefore be enough to prove that $\PP(M < \infty) > 0$, or equivalently that $\PP(M = \infty) < 1$. In the remaining part of the proof we verify that $\PP(M = \infty) < 1$.

For every $k \in \NN$, set
\[ M_k := \sup_{0 \le n \le k} (S_n - na), \qquad M_k' := \sup_{0 \le n \le k} (S_{n+1} - S_1 - na). \]
Then, $M_k$ and $M_k'$ are nonnegative random variables with the same distribution, as a consequence of the fact that the random vectors $(X_1,\ldots,X_k)$ and $(X_2,\ldots,X_{k+1})$ have the same law. It follows that
\[ M = \lim_{k \to \infty} \uparrow\, M_k \qquad \text{and} \qquad M' := \lim_{k \to \infty} \uparrow\, M_k' \]
also have the same law (just observe that $\PP(M' \le x) = \lim \downarrow \PP(M_k' \le x) = \lim \downarrow \PP(M_k \le x) = \PP(M \le x)$ for every $x \in \RR$).

On the other hand, it follows from the definitions that, for every integer $k \ge 0$,
\[ M_{k+1} = \sup\left(0,\ \sup_{1 \le n \le k+1} (S_n - na)\right) = \sup(0, M_k' + X_1 - a), \]
which can also be written as
\[ M_{k+1} = M_k' - \inf(a - X_1, M_k'). \]
Since $M_k'$ has the same law as $M_k$ (and both these variables are clearly in $L^1$), we get
\[ \EE[\inf(a - X_1, M_k')] = \EE[M_k'] - \EE[M_{k+1}] = \EE[M_k] - \EE[M_{k+1}] \le 0 \]
thanks to the trivial inequality $M_k \le M_{k+1}$. We may now apply the dominated convergence theorem to the sequence $\inf(a - X_1, M_k')$, $k \in \NN$, noting that these variables are bounded in absolute value by $|a - X_1|$ (recall that $M_k' \ge 0$). It follows that
\[ \EE[\inf(a - X_1, M')] = \lim_{k \to \infty} \EE[\inf(a - X_1, M_k')] \le 0. \]

We now argue by contradiction to prove that $\PP(M = \infty) < 1$. Suppose that $\PP(M = \infty) = 1$, then we have also $\PP(M' = \infty) = 1$, since $M$ and $M'$ have the same law, and it follows that $\inf(a - X_1, M') = a - X_1$ a.s. But then the inequality in the last display gives $\EE[a - X_1] \le 0$, which is a contradiction since we chose $a > \EE[X_1]$. This contradiction completes the proof.
}

\textbf{Remarks.}
\begin{itemize}
    \item  The $L^1$ integrability assumption is optimal in the sense that it is needed for the limit $\EE[X_1]$ to be defined (and finite). If we remove the $L^1$ assumption but assume instead that the random variables $X_n$ are nonnegative and $\EE[X_1] = \infty$, it is easy to obtain that
    \[ \frac{1}{n}(X_1 + \cdots + X_n) \xrightarrow[n \to \infty]{\mathrm{a.s.}} +\infty, \]
    by applying the theorem to the variables $X_n \wedge K$, for every $K \in \NN$.

    \item  One can show that the convergence of the theorem also holds in $L^1$. The proof will be given at the end of next chapter as an application of martingale theory. As we already mentioned, the almost sure convergence is the most interesting one from a probabilistic point of view.
\end{itemize}



\subsection{The Central Limit Theorem}
Let $(X_n)_{n \in \NN}$ be a sequence of independent and identically distributed real random variables in $L^1$. The strong law of large numbers asserts that
\[ \frac{1}{n}(X_1 + \cdots + X_n) \xrightarrow[n \to \infty]{a.s.} \EE[X_1]. \]

We are now interested in the speed of this convergence, meaning that we want information about the typical size of
\[
\frac{1}{n}(X_1 + \cdots + X_n) - \EE[X_1]
\]
when $n$ is large.

Under the additional assumption that the variables $X_n$ are in $L^2$, it is easy to guess the answer, because with a simple calculation we can conclude that 
\[
\EE\left[(X_1 + \cdots + X_n - n\EE[X_1])^2\right] = \mathrm{var}(X_1 + \cdots + X_n) = n\,\mathrm{var}(X_1).
\]

This shows that the mean value of $(X_1 + \cdots + X_n - n\EE[X_1])^2$ grows linearly with $n$, and thus the typical size of $X_1 + \cdots + X_n - n\EE[X_1]$ is expected to be $\sqrt{n}$, or equivalently the typical size of
\[ \frac{1}{n}(X_1 + \cdots + X_n) - \EE[X_1] \]
should be $1/\sqrt{n}$. The central limit theorem gives a much more precise statement.

\thmpr{Central Limit Theorem}{central_limit_thm}{
Let $(X_n)_{n \in \NN}$ be a sequence of independent and identically distributed real random variables in $L^2$. Set $\sigma^2 = \mathrm{var}(X_1)$ and assume that $\sigma > 0$. Then
\[ \frac{1}{\sqrt{n}}(X_1 + \cdots + X_n - n\EE[X_1]) \conv[\D]{n} \N(0,\sigma^2). \]

where $\N(0,\sigma^2)$ is the Gaussian distribution with mean $0$ and variance $\sigma^2$. Hence, for every $a,b \in \RR$ with $a < b$,
\[ \lim_{n \to \infty} \PP(X_1 + \cdots + X_n \in [n\EE[X_1] + a\sqrt{n},\, n\EE[X_1] + b\sqrt{n}]) = \frac{1}{\sigma\sqrt{2\pi}} \int_a^b \exp\left(-\frac{x^2}{2\sigma^2}\right)\,dx. \]
}{
The second part of the statement follows from the first one, thanks to Theorem \ref{thm:Portmanteau_theorem}. To prove the first assertion, note that we can assume $\EE[X_1] = 0$ (just replace $X_n$ by $X_n - \EE[X_n]$). Then set
\[ Z_n = \frac{1}{\sqrt{n}}(X_1 + \cdots + X_n). \]

The characteristic function of $Z_n$ is
\begin{align*}
    \Phi_{Z_n}(\xi) &= \EE\left[\exp\left(i\xi \frac{X_1 + \cdots + X_n}{\sqrt{n}}\right)\right] \\
    &= \EE\left[\exp\left(i\frac{\xi}{\sqrt{n}} X_1\right)\right]^n \\
    &= \Phi_{X_1}\left(\frac{\xi}{\sqrt{n}}\right)^n.
\end{align*}

where, in the second equality, we used the fact that the random variables $X_i$ are independent and identically distributed. On the other hand, by Proposition \ref{prop:taylor_expansion_characteristic_function}, we have
\[ \Phi_{X_1}(\xi) = 1 + i\xi \EE[X_1] - \frac{1}{2}\xi^2 \EE[X_1^2] + o(\xi^2) = 1 - \frac{\sigma^2 \xi^2}{2} + o(\xi^2) \]
as $\xi \to 0$. For every fixed $\xi \in \RR$, we have thus
\[ \Phi_{X_1}\left(\frac{\xi}{\sqrt{n}}\right) = 1 - \frac{\sigma^2 \xi^2}{2n} + o\left(\frac{1}{n}\right) \]
as $n \to \infty$. Thanks to proposition \ref{prop:prob_imply_Lp} and the last display, we have, for every $\xi \in \RR$,
\[ \lim_{n \to \infty} \Phi_{Z_n}(\xi) = \lim_{n \to \infty} \left(1 - \frac{\sigma^2 \xi^2}{2n} + o\left(\frac{1}{n}\right)\right)^n = \exp\left(-\frac{\sigma^2 \xi^2}{2}\right), \]
which is the characteristic function of the Gaussian $\N(0,\sigma^2)$ distribution as described in appendix \ref{sec:gaussian_sec}. An application of Lévy's theorem \ref{thm:Levy_theorem} now completes the proof.
}
\textbf{Remark} When $\sigma = 0$ there is not much to say since the variables $X_n$ are all equal a.s. to the constant $\EE[X_1]$.


\subsection{The Multidimensional Central Limit Theorem}
We now assume that $(X_n)_{n \in \NN}$ is a sequence of independent and identically distributed random variables with values in $\RR^d$, such that $\EE[\|X_1\|] < \infty$. We can apply the strong law of large numbers to each component of $X_n$ and we have
\[
\frac{1}{n}(X_1 + \cdots + X_n) \xrightarrow[n \to \infty]{a.s.} \EE[X_1],
\]
as in the real case. If we assume furthermore that $\EE[\|X_1\|^2] < \infty$, then we expect to have also an analog of the central limit theorem. However, in contrast with the law of large numbers, this is not immediate because the convergence in distribution of a sequence of random vectors does not follow from the convergence in distribution of each component sequence.

To extend the central limit theorem to the case of random variables with values in $\RR^d$, we recall the main 3 results of multidimensional Gaussian distributions :
\begin{enumerate}
    \item A random variable $\bX=(X_1,\ldots,X_p)$ has a multidimensional Gaussian distribution (ie Gaussian vector) iff for all $a_1,\ldots,a_p\in\RR$ the real random variable $a_1X_1+\ldots+a_pX_p$ has a Gaussian distribution.
    \item The characteristic function of Gaussian vector $\N(m_\bX, K_\bX)$ is provided by :
    \[ \Phi_{\bX}(\xi) = \EE\left[ e^{\xi\cdot \bX} \right] = e^{i\xi\cdot m_\bX -\frac{1}{2} {}^t\xi K_\bX \xi}   \] 
    \item If $\Gamma$ be a symmetric nonnegative definite $d\times d$ matrix. Then there exists a Gaussian vector $\bX \sim \N_d(0,\Gamma)$.
\end{enumerate}


\thmpr{Multidimensional Central Limit Theorem}{ndim_central_limit_thm}{
Let $(X_n)_{n \in \NN}$ be a sequence of independent and identically distributed random variables with values in $\RR^d$, such that $\EE[\|X_1\|^2] < \infty$. Then,
\[ \frac{1}{\sqrt{n}}(X_1 + \cdots + X_n - n\EE[X_1]) \conv[\D]{n} \N_d(0, K_{X_1}). \]
}{
Replacing $X_n$ by $X_n - \EE[X_n]$ allows us to assume that $\EE[X_1] = 0$. For every $\xi \in \RR^d$, the central limit theorem in the real case (Theorem \ref{thm:central_limit_thm}) shows that
\[
\frac{1}{\sqrt{n}}(\xi \cdot X_1 + \cdots + \xi \cdot X_n)
\conv[\D]{n} \N(0,\sigma^2),
\]
where $\sigma^2 = \EE[(\xi \cdot X_1)^2] = {}^t\xi K_{X_1} \xi$. This implies that
\[
\EE\left[\exp\left(i\xi \cdot \frac{X_1 + \cdots + X_n}{\sqrt{n}}\right)\right] \conv{n} \exp\left(-\frac{1}{2}\,{}^t\xi K_{X_1}\xi\right)
\]
and Lévy's theorem (Theorem \ref{thm:Levy_theorem}) gives the desired result.
}




\end{document}